\documentclass[8pt]{beamer}
\usetheme{Madrid}
\usecolortheme{default}
\usepackage{amsmath, amssymb, graphicx, array, multirow, booktabs, xcolor, tikz}
\usetikzlibrary{shapes,arrows,positioning,fit,backgrounds}

% Compact font settings
\setbeamerfont{title}{size=\Large}
\setbeamerfont{subtitle}{size=\small}
\setbeamerfont{author}{size=\scriptsize}
\setbeamerfont{date}{size=\scriptsize}
\setbeamerfont{frametitle}{size=\large}
\setbeamerfont{framesubtitle}{size=\normalsize}
\setbeamerfont{enumerate item}{size=\footnotesize}
\setbeamerfont{itemize item}{size=\footnotesize}
\setbeamerfont{block title}{size=\footnotesize}

% Compact spacing
\setlength{\itemsep}{0.08ex}
\setlength{\parskip}{0.08ex}
\setlength{\leftmargini}{0.3cm}
\setbeamersize{text margin left=3mm, text margin right=3mm}
\addtobeamertemplate{frame begin}{\vspace{-0.4cm}}{}

% Colors
\definecolor{myblue}{RGB}{0,0,255}
\definecolor{myred}{RGB}{255,0,0}
\definecolor{mygreen}{RGB}{0,128,0}
\definecolor{myorange}{RGB}{255,165,0}
\definecolor{mypurple}{RGB}{128,0,128}
\definecolor{mygray}{RGB}{128,128,128}
\definecolor{mygold}{RGB}{255,215,0}

\title{Chapter 7: Supervised Learning - Model Estimation}
\subtitle{100 Tutorial Topics: OLS, NLS, MLE, Gradient Descent, Backpropagation, Regularization, Cross-Validation}
\author{GARP RAI Exam Preparation}
\date{}

\begin{document}
	
	\begin{frame}
		\titlepage
	\end{frame}
	
	% ============================================================
	% SECTION 1: INTRODUCTION TO MODEL ESTIMATION
	% ============================================================
	
	\begin{frame}{T1: Model Estimation Purpose}
		\fontsize{6.5}{7.5}\selectfont
		\textbf{TOPIC: Primary Purpose of Model Estimation in Supervised Learning}
		
		\vspace{0.1cm}
		\textbf{DEFINITION:}
		\begin{itemize}
			\item Model estimation is the essential step where a model learns from data
			\item It involves finding the specific parameter values that cause the model to best fit the training data
			\item Parameters are often called "weights" in machine learning
		\end{itemize}
		
		\vspace{0.1cm}
		\textbf{TWO MAIN ESTIMATION FAMILIES:}
		\begin{enumerate}
			\item \textcolor{myblue}{Least Squares:} Minimize sum of squared differences between predictions and actual values
			\begin{itemize}
				\item Objective: Make prediction errors (residuals) as small as possible
			\end{itemize}
			\item \textcolor{myblue}{Maximum Likelihood:} Probabilistic approach
			\begin{itemize}
				\item Assume data comes from specific distribution
				\item Find parameters that make observing data most likely
			\end{itemize}
		\end{enumerate}
		
		\vspace{0.1cm}
		\textbf{WHY IT MATTERS:}
		\begin{itemize}
			\item Fundamental to all supervised learning
			\item Determines model performance
			\item Foundation for prediction and classification
		\end{itemize}
		
		\textcolor{myred}{\textbf{EXAM TRAP:}} Model estimation is about finding \textcolor{myblue}{optimal parameter values}, not increasing complexity!
	\end{frame}
	
	\begin{frame}{T2: Analytical vs Numerical Methods}
		\fontsize{6.5}{7.5}\selectfont
		\textbf{TOPIC: Distinguishing Between Analytical and Numerical Estimation Methods}
		
		\vspace{0.1cm}
		\textbf{DEFINITION:}
		\begin{itemize}
			\item Estimation methods can be broadly classified into analytical and numerical approaches
			\item The choice depends on the complexity of the objective function and model
		\end{itemize}
		
		\vspace{0.1cm}
		\textbf{ANALYTICAL METHODS:}
		\begin{itemize}
			\item \textbf{Closed-form solution:} Exact mathematical equations
			\item \textbf{Unique solution:} Parameter values are unique
			\item \textbf{Example:} OLS for linear regression: $\hat{\beta} = (X^TX)^{-1}X^Ty$
			\item \textbf{When available:} When objective function can be differentiated
		\end{itemize}
		
		\vspace{0.1cm}
		\textbf{NUMERICAL METHODS:}
		\begin{itemize}
			\item \textbf{No closed-form solution:} Iterative procedures
			\item \textbf{Start with initial guess:} Gradually move toward optimal
			\item \textbf{Example:} Gradient descent, NLS
			\item \textbf{When needed:} When differentiation is too complex or impossible
		\end{itemize}
		
		\vspace{0.1cm}
		\textbf{KEY INSIGHT:} Analytical = exact formula; Numerical = iterative approximation
		
		\textcolor{myred}{\textbf{EXAM TRAP:}} Analytical methods have \textcolor{myblue}{closed-form solutions}; numerical methods require \textcolor{myblue}{iterative procedures}!
	\end{frame}
	
	\begin{frame}{T3: Hyperparameters vs Parameters}
		\fontsize{6.5}{7.5}\selectfont
		\textbf{TOPIC: Distinguishing Between Hyperparameters and Parameters}
		
		\vspace{0.1cm}
		\textbf{DEFINITION:}
		\begin{itemize}
			\item Parameters are learned from data during training
			\item Hyperparameters are configuration settings chosen before training begins
		\end{itemize}
		
		\vspace{0.1cm}
		\textbf{PARAMETERS (WEIGHTS):}
		\begin{itemize}
			\item \textbf{Learned from data:} Estimated by optimization algorithm
			\item \textbf{Examples:}
			\begin{itemize}
				\item Regression coefficients ($\beta_0, \beta_1, ...$)
				\item Neural network weights
				\item Bias terms
			\end{itemize}
			\item \textbf{Updated during training:} Changed with each iteration
		\end{itemize}
		
		\vspace{0.1cm}
		\textbf{HYPERPARAMETERS:}
		\begin{itemize}
			\item \textbf{Set before training:} Configuration settings
			\item \textbf{Examples:}
			\begin{itemize}
				\item Learning rate ($\eta$)
				\item Number of hidden layers
				\item Regularization parameter ($\lambda$)
				\item Number of epochs
			\end{itemize}
			\item \textbf{Require tuning:} Grid search, cross-validation
		\end{itemize}
		
		\textcolor{myred}{\textbf{EXAM TRAP:}} Parameters = \textcolor{myblue}{learned from data}; Hyperparameters = \textcolor{myblue}{set before training}!
	\end{frame}
	
	\begin{frame}{T4: Least Squares Overview}
		\fontsize{6.5}{7.5}\selectfont
		\textbf{TOPIC: Overview of the Least Squares Method}
		
		\vspace{0.1cm}
		\textbf{DEFINITION:}
		\begin{itemize}
			\item Technique to estimate parameters by minimizing sum of squared differences between observed values and model predictions
			\item Two variants: Ordinary Least Squares (OLS) and Nonlinear Least Squares (NLS)
		\end{itemize}
		
		\vspace{0.1cm}
		\textbf{OBJECTIVE FUNCTION:}
		\begin{itemize}
			\item Residual Sum of Squares (RSS): $\sum_{i=1}^N (y_i - \hat{y}_i)^2$
			\item Mean Squared Error (MSE): $\frac{1}{N}\sum_{i=1}^N (y_i - \hat{y}_i)^2$
		\end{itemize}
		
		\vspace{0.1cm}
		\textbf{WHY SQUARE RESIDUALS?}
		\begin{itemize}
			\item Positive and negative residuals would cancel if simply summed
			\item Sum of residuals is always zero at optimal solution
			\item Squaring ensures all residuals contribute positively
			\item Penalizes large errors more heavily
		\end{itemize}
		
		\vspace{0.1cm}
		\textbf{KEY CHARACTERISTICS:}
		\begin{itemize}
			\item \textcolor{myblue}{OLS:} Linear in parameters; closed-form solution
			\item \textcolor{myblue}{NLS:} Nonlinear in parameters; iterative solution
		\end{itemize}
		
		\textcolor{myred}{\textbf{EXAM TRAP:}} Least Squares minimizes \textcolor{myblue}{squared} differences, not absolute differences!
	\end{frame}
	
	\begin{frame}{T5: Maximum Likelihood Overview}
		\fontsize{6.5}{7.5}\selectfont
		\textbf{TOPIC: Overview of Maximum Likelihood Estimation}
		
		\vspace{0.1cm}
		\textbf{DEFINITION:}
		\begin{itemize}
			\item Probabilistic approach to parameter estimation
			\item Finds parameter values that maximize probability of observing the given data
			\item More flexible framework than Least Squares
		\end{itemize}
		
		\vspace{0.1cm}
		\textbf{CORE QUESTION:}
		\begin{itemize}
			\item "Given observed data, what parameter values make this data most probable?"
		\end{itemize}
		
		\vspace{0.1cm}
		\textbf{PROCESS:}
		\begin{enumerate}
			\item Assume data comes from specific distribution
			\item Construct likelihood function: $L(\theta) = P(X|\theta)$
			\item Find $\theta$ that maximizes $L(\theta)$
			\item Use log-likelihood for easier computation
		\end{enumerate}
		
		\vspace{0.1cm}
		\textbf{APPLICATIONS:}
		\begin{itemize}
			\item Logistic regression (Bernoulli distribution)
			\item Linear regression (Normal distribution)
			\item Probit models
			\item GARCH models for volatility
		\end{itemize}
		
		\textcolor{myred}{\textbf{EXAM TRAP:}} MLE maximizes \textcolor{myblue}{likelihood} = $P(\text{data}|\theta)$, NOT $P(\theta|\text{data})$!
	\end{frame}
	
	% ============================================================
	% SECTION 2: ORDINARY LEAST SQUARES
	% ============================================================
	
	\begin{frame}{T6: OLS Definition}
		\fontsize{6.5}{7.5}\selectfont
		\textbf{TOPIC: Ordinary Least Squares - Definition and Overview}
		
		\vspace{0.1cm}
		\textbf{DEFINITION:}
		\begin{itemize}
			\item Most straightforward estimation method
			\item Analytical approach with closed-form solution
			\item Used to estimate parameters in linear regression models
		\end{itemize}
		
		\vspace{0.1cm}
		\textbf{OBJECTIVE:}
		\begin{itemize}
			\item Minimize Residual Sum of Squares (RSS)
			\item $RSS = \sum_{i=1}^N (y_i - \hat{y}_i)^2 = \sum_{i=1}^N (y_i - \hat{\beta}_0 - \hat{\beta}_1 x_i)^2$
		\end{itemize}
		
		\vspace{0.1cm}
		\textbf{KEY CHARACTERISTICS:}
		\begin{itemize}
			\item \textcolor{myblue}{Analytical method:} Closed-form solution
			\item \textcolor{myblue}{Linear in parameters:} $y = \beta_0 + \beta_1 x + \varepsilon$
			\item \textcolor{myblue}{Unique solution:} One optimal set of parameters
			\item \textcolor{myblue}{Computationally efficient:} Direct calculation
		\end{itemize}
		
		\vspace{0.1cm}
		\textbf{EXAMPLE:}
		\begin{itemize}
			\item Salary = 16.88 + 1.06 × Experience
			\item Intercept = 16.88, Slope = 1.06
		\end{itemize}
		
		\textcolor{myred}{\textbf{EXAM TRAP:}} OLS is an \textcolor{myblue}{analytical} method with \textcolor{myblue}{closed-form solution}!
	\end{frame}
	
	\begin{frame}{T7: OLS - Residual Sum of Squares}
		\fontsize{6.5}{7.5}\selectfont
		\textbf{TOPIC: Residual Sum of Squares (RSS) in OLS}
		
		\vspace{0.1cm}
		\textbf{DEFINITION:}
		\begin{itemize}
			\item RSS is the sum of squared differences between actual and predicted values
			\item It is the objective function that OLS minimizes
		\end{itemize}
		
		\vspace{0.1cm}
		\textbf{FORMULA:}
		\[
		RSS = \sum_{i=1}^N \hat{u}_i^2 = \sum_{i=1}^N (y_i - \hat{y}_i)^2
		\]
		where:
		\begin{itemize}
			\item $y_i$ = actual observed value
			\item $\hat{y}_i$ = predicted value
			\item $\hat{u}_i = y_i - \hat{y}_i$ = residual
		\end{itemize}
		
		\vspace{0.1cm}
		\textbf{RELATED METRICS:}
		\begin{itemize}
			\item MSE = RSS/N (Mean Squared Error)
			\item RMSE = $\sqrt{MSE}$ (Root Mean Squared Error)
		\end{itemize}
		
		\vspace{0.1cm}
		\textbf{WHY SQUARE?}
		\begin{itemize}
			\item Prevents positive and negative residuals from canceling
			\item Sum of residuals is always zero at OLS solution
			\item Penalizes large errors more heavily
		\end{itemize}
		
		\textcolor{myred}{\textbf{EXAM TRAP:}} RSS = \textcolor{myblue}{sum of squared residuals}, NOT sum of absolute residuals!
	\end{frame}
	
	\begin{frame}{T8: OLS - MSE Formula}
		\fontsize{6.5}{7.5}\selectfont
		\textbf{TOPIC: Mean Squared Error (MSE) in OLS}
		
		\vspace{0.1cm}
		\textbf{DEFINITION:}
		\begin{itemize}
			\item MSE is the average of squared residuals
			\item Commonly used metric for model performance
		\end{itemize}
		
		\vspace{0.1cm}
		\textbf{FORMULA:}
		\[
		MSE = \frac{1}{N}\sum_{i=1}^N \hat{u}_i^2 = \frac{1}{N}\sum_{i=1}^N (y_i - \hat{y}_i)^2 = \frac{RSS}{N}
		\]
		
		\vspace{0.1cm}
		\textbf{RELATIONSHIP TO RSS:}
		\begin{itemize}
			\item MSE = RSS / N
			\item Both minimized at same parameter values
			\item MSE is scale-independent of sample size
		\end{itemize}
		
		\vspace{0.1cm}
		\textbf{WHY MSE IS USEFUL:}
		\begin{itemize}
			\item Easier to compare across datasets
			\item Commonly used as loss function in machine learning
			\item Basis for RMSE calculation
		\end{itemize}
		
		\textcolor{myred}{\textbf{EXAM TRAP:}} MSE = \textcolor{myblue}{RSS/N} = average squared error!
	\end{frame}
	
	\begin{frame}{T9: OLS - Fitted Line Example}
		\fontsize{6.5}{7.5}\selectfont
		\textbf{TOPIC: Fitted Regression Line in Salary Example}
		
		\vspace{0.1cm}
		\textbf{EXAMPLE DATA:}
		\begin{itemize}
			\item 14 observations of bank employee experience and salary
			\item Simple linear regression: one feature (experience), one target (salary)
		\end{itemize}
		
		\vspace{0.1cm}
		\textbf{FITTED EQUATION:}
		\[
		\hat{y} = 16.88 + 1.06x
		\]
		where:
		\begin{itemize}
			\item $\hat{y}$ = predicted salary (USD/hour)
			\item $x$ = years of experience
			\item Intercept $\hat{\beta}_0 = 16.88$
			\item Slope $\hat{\beta}_1 = 1.06$
		\end{itemize}
		
		\vspace{0.1cm}
		\textbf{INTERPRETATION:}
		\begin{itemize}
			\item Base salary (0 experience): \$16.88/hour
			\item Each additional year: +\$1.06/hour
		\end{itemize}
		
		\vspace{0.1cm}
		\textbf{EXAMPLE PREDICTION:}
		\begin{itemize}
			\item Experience = 9 years
			\item Predicted salary = 16.88 + 9(1.06) = 26.42
		\end{itemize}
		
		\textcolor{myred}{\textbf{EXAM TRAP:}} Know the \textcolor{myblue}{intercept and slope} from the example!
	\end{frame}
	
	\begin{frame}{T10: OLS - Residual Calculation}
		\fontsize{6.5}{7.5}\selectfont
		\textbf{TOPIC: Calculating Residuals in OLS}
		
		\vspace{0.1cm}
		\textbf{DEFINITION:}
		\begin{itemize}
			\item Residual = Actual value - Predicted value
			\item $\hat{u}_i = y_i - \hat{y}_i$
		\end{itemize}
		
		\vspace{0.1cm}
		\textbf{EXAMPLE:}
		\begin{itemize}
			\item Data point: Experience = 9 years, Salary = \$33.08/hour
			\item Predicted: $\hat{y} = 16.88 + 9(1.06) = 26.42$
			\item Residual: $33.08 - 26.42 = 6.66$
		\end{itemize}
		
		\vspace{0.1cm}
		\textbf{INTERPRETATION:}
		\begin{itemize}
			\item Positive residual: actual > predicted
			\item Model under-predicted for this point
			\item Negative residual: actual < predicted
			\item Model over-predicted for this point
		\end{itemize}
		
		\vspace{0.1cm}
		\textbf{PROPERTY:}
		\begin{itemize}
			\item Sum of residuals = 0 at OLS solution
			\item $\sum \hat{u}_i = 0$
		\end{itemize}
		
		\textcolor{myred}{\textbf{EXAM TRAP:}} Residual = \textcolor{myblue}{actual - predicted}, NOT predicted - actual!
	\end{frame}
	
	\begin{frame}{T11: OLS - Normal Equations}
		\fontsize{6.5}{7.5}\selectfont
		\textbf{TOPIC: Deriving OLS Estimates Using Normal Equations}
		
		\vspace{0.1cm}
		\textbf{DERIVATION PROCESS:}
		\begin{enumerate}
			\item Write RSS as function of parameters
			\item Partially differentiate with respect to each parameter
			\item Set derivatives to zero (first-order conditions)
			\item Solve system of equations (normal equations)
		\end{enumerate}
		
		\vspace{0.1cm}
		\textbf{RESULTING FORMULAS:}
		\begin{itemize}
			\item Intercept: $\hat{\beta}_0 = \bar{y} - \hat{\beta}_1 \bar{x}$
			\item Slope: $\hat{\beta}_1 = \frac{\sum_{i=1}^N y_i x_i - N\bar{y}\bar{x}}{\sum_{i=1}^N x_i^2 - N\bar{x}^2}$
		\end{itemize}
		
		\vspace{0.1cm}
		\textbf{WHY THIS WORKS:}
		\begin{itemize}
			\item RSS is convex function of parameters
			\item Unique minimum exists
			\item First-order conditions characterize minimum
		\end{itemize}
		
		\vspace{0.1cm}
		\textbf{KEY INSIGHT:}
		\begin{itemize}
			\item OLS has closed-form solution
			\item No iterative optimization needed
			\item Direct calculation from data
		\end{itemize}
		
		\textcolor{myred}{\textbf{EXAM TRAP:}} OLS uses \textcolor{myblue}{calculus} (derivatives) to find closed-form solution!
	\end{frame}
	
	\begin{frame}{T12: OLS - Multiple Regression}
		\fontsize{6.5}{7.5}\selectfont
		\textbf{TOPIC: Extending OLS to Multiple Regression}
		
		\vspace{0.1cm}
		\textbf{DEFINITION:}
		\begin{itemize}
			\item Multiple regression has more than one explanatory variable
			\item Model: $y = \beta_0 + \beta_1 x_1 + \beta_2 x_2 + ... + \beta_m x_m + \varepsilon$
		\end{itemize}
		
		\vspace{0.1cm}
		\textbf{MATRIX FORMULATION:}
		\begin{itemize}
			\item Model: $y = X\beta + \varepsilon$
			\item Normal equations: $X^TX\hat{\beta} = X^Ty$
			\item Solution: $\hat{\beta} = (X^TX)^{-1}X^Ty$
		\end{itemize}
		
		\vspace{0.1cm}
		\textbf{WHY MATRIX ALGEBRA:}
		\begin{itemize}
			\item Formulas become complex with many variables
			\item Matrix notation is compact
			\item Computationally efficient
		\end{itemize}
		
		\vspace{0.1cm}
		\textbf{REQUIREMENTS:}
		\begin{itemize}
			\item $X^TX$ must be invertible
			\item No perfect multicollinearity
		\end{itemize}
		
		\textcolor{myred}{\textbf{EXAM TRAP:}} Multiple regression uses \textcolor{myblue}{matrix algebra}, not separate regressions!
	\end{frame}
	
	\begin{frame}{T13: OLS - Goodness of Fit}
		\fontsize{6.5}{7.5}\selectfont
		\textbf{TOPIC: Interpreting OLS Goodness of Fit}
		
		\vspace{0.1cm}
		\textbf{RSS INTERPRETATION:}
		\begin{itemize}
			\item \textcolor{myblue}{Low RSS:} Small residuals, good fit
			\item \textcolor{myblue}{High RSS:} Large residuals, poor fit
			\item \textcolor{myblue}{RSS = 0:} Perfect fit (rare)
		\end{itemize}
		
		\vspace{0.1cm}
		\textbf{CAVEATS:}
		\begin{itemize}
			\item Low training RSS could indicate overfitting
			\item Need to check test/validation performance
			\item Compare to baseline models
		\end{itemize}
		
		\vspace{0.1cm}
		\textbf{SALARY EXAMPLE:}
		\begin{itemize}
			\item Large residuals observed
			\item Suggested OLS not ideal fit
			\item Motivated considering alternative models
		\end{itemize}
		
		\vspace{0.1cm}
		\textbf{RELATED METRICS:}
		\begin{itemize}
			\item R-squared: Proportion of variance explained
			\item Adjusted R-squared: Penalizes additional parameters
		\end{itemize}
		
		\textcolor{myred}{\textbf{EXAM TRAP:}} Low RSS = \textcolor{myblue}{good fit}, but check for overfitting!
	\end{frame}
	
	\begin{frame}{T14: OLS - Limitations}
		\fontsize{6.5}{7.5}\selectfont
		\textbf{TOPIC: Key Limitations of OLS}
		
		\vspace{0.1cm}
		\textbf{LINEARITY REQUIREMENT:}
		\begin{itemize}
			\item OLS requires model linear in parameters
			\item $y = \beta_0 + \beta_1 x + \beta_2 x^2$ is OK (linear in $\beta$s)
			\item $y = \beta_0 + \beta_1 e^{\beta_2 x}$ is NOT OK (nonlinear in $\beta_2$)
		\end{itemize}
		
		\vspace{0.1cm}
		\textbf{OTHER LIMITATIONS:}
		\begin{itemize}
			\item Sensitive to outliers
			\item Assumes homoscedasticity (constant error variance)
			\item Assumes no autocorrelation
			\item No closed-form for nonlinear models
		\end{itemize}
		
		\vspace{0.1cm}
		\textbf{SOLUTION FOR NONLINEAR:}
		\begin{itemize}
			\item Use Nonlinear Least Squares (NLS)
			\item Use numerical optimization
			\item Use Maximum Likelihood Estimation
		\end{itemize}
		
		\vspace{0.1cm}
		\textbf{BINARY OUTCOMES:}
		\begin{itemize}
			\item OLS cannot handle binary targets
			\item Use logistic regression (MLE)
		\end{itemize}
		
		\textcolor{myred}{\textbf{EXAM TRAP:}} OLS requires \textcolor{myblue}{linearity in parameters}, not variables!
	\end{frame}
	
	\begin{frame}{T15: OLS vs NLS}
		\fontsize{6.5}{7.5}\selectfont
		\textbf{TOPIC: Comparing OLS and NLS}
		
		\vspace{0.1cm}
		\textbf{KEY DIFFERENCES:}
		\begin{tabular}{|p{2.5cm}|p{3cm}|p{3cm}|}
			\hline
			\textbf{Aspect} & \textbf{OLS} & \textbf{NLS} \\
			\hline
			Linearity & Linear in parameters & Nonlinear in parameters \\
			\hline
			Solution & Closed-form & Iterative \\
			\hline
			Method & Analytical & Numerical \\
			\hline
			Speed & Fast & Slower \\
			\hline
			Example & $y = \beta_0 + \beta_1 x$ & $y = \beta_0 + \beta_1 e^{\beta_2 x}$ \\
			\hline
		\end{tabular}
		
		\vspace{0.1cm}
		\textbf{COMMON GROUND:}
		\begin{itemize}
			\item Both minimize sum of squared errors
			\item Both are least squares methods
		\end{itemize}
		
		\vspace{0.1cm}
		\textbf{WHEN TO USE EACH:}
		\begin{itemize}
			\item \textcolor{myblue}{OLS:} Simple linear regression, polynomial regression
			\item \textcolor{myblue}{NLS:} Neural networks, exponential models
		\end{itemize}
		
		\textcolor{myred}{\textbf{EXAM TRAP:}} OLS = \textcolor{myblue}{closed-form}; NLS = \textcolor{myblue}{iterative}!
	\end{frame}
	
	\begin{frame}{T16: NLS Definition}
		\fontsize{6.5}{7.5}\selectfont
		\textbf{TOPIC: Nonlinear Least Squares - Definition and Process}
		
		\vspace{0.1cm}
		\textbf{DEFINITION:}
		\begin{itemize}
			\item Iterative numerical approach to estimate parameters in nonlinear models
			\item Minimizes objective function (RSS or MSE)
			\item No closed-form solution exists
		\end{itemize}
		
		\vspace{0.1cm}
		\textbf{NLS PROCESS:}
		\begin{enumerate}
			\item Start with initial "guesses" for parameters (small random numbers)
			\item Make predictions and evaluate objective function
			\item Adjust parameters to improve objective
			\item Check convergence; repeat if needed
		\end{enumerate}
		
		\vspace{0.1cm}
		\textbf{WHY ITERATIVE:}
		\begin{itemize}
			\item No closed-form solution
			\item Normal equations are nonlinear
			\item Requires numerical optimization
		\end{itemize}
		
		\vspace{0.1cm}
		\textbf{APPLICATIONS:}
		\begin{itemize}
			\item Neural networks
			\item Complex curve fitting
			\item Models with exponential terms
		\end{itemize}
		
		\textcolor{myred}{\textbf{EXAM TRAP:}} NLS is \textcolor{myblue}{iterative} because no closed-form solution exists!
	\end{frame}
	
	\begin{frame}{T17: NLS - Hill Climbing}
		\fontsize{6.5}{7.5}\selectfont
		\textbf{TOPIC: Hill Climbing in NLS Optimization}
		
		\vspace{0.1cm}
		\textbf{DEFINITION:}
		\begin{itemize}
			\item Simple method for estimating parameters
			\item Tests small changes to each parameter one at a time
			\item No derivatives required
		\end{itemize}
		
		\vspace{0.1cm}
		\textbf{APPROACH:}
		\begin{enumerate}
			\item Start with initial guesses
			\item Test small changes to each parameter
			\item Keep changes that improve objective
			\item Repeat until no improvement
		\end{enumerate}
		
		\vspace{0.1cm}
		\textbf{ADVANTAGES:}
		\begin{itemize}
			\item Simple to implement
			\item No derivatives needed
			\item Works for non-differentiable functions
		\end{itemize}
		
		\vspace{0.1cm}
		\textbf{DISADVANTAGES:}
		\begin{itemize}
			\item Gets stuck in local optima
			\item Slow convergence
			\item Only one parameter at a time
			\item Misses optimal combinations
		\end{itemize}
		
		\textcolor{myred}{\textbf{EXAM TRAP:}} Hill climbing is \textcolor{myblue}{simple but prone to local optima}!
	\end{frame}
	
	\begin{frame}{T18: NLS - Gradient Descent Advantage}
		\fontsize{6.5}{7.5}\selectfont
		\textbf{TOPIC: Why Gradient Descent is Preferred Over Hill Climbing}
		
		\vspace{0.1cm}
		\textbf{GRADIENT DESCENT ADVANTAGES:}
		\begin{itemize}
			\item Updates all parameters simultaneously
			\item Uses gradient information (steepest descent)
			\item Faster convergence than hill climbing
			\item More efficient for complex models
		\end{itemize}
		
		\vspace{0.1cm}
		\textbf{HILL CLIMBING LIMITATIONS:}
		\begin{itemize}
			\item Only one parameter at a time
			\item Misses optimal parameter combinations
			\item Slower convergence
		\end{itemize}
		
		\vspace{0.1cm}
		\textbf{COMMON GROUND:}
		\begin{itemize}
			\item Both can get stuck in local minima
			\item Neither guarantees global minimum
		\end{itemize}
		
		\vspace{0.1cm}
		\textbf{WHY GRADIENT DESCENT WINS:}
		\begin{itemize}
			\item Workhorse of modern machine learning
			\item Scales to high dimensions
			\item Foundation of neural network training
		\end{itemize}
		
		\textcolor{myred}{\textbf{EXAM TRAP:}} Gradient descent updates \textcolor{myblue}{all parameters simultaneously}!
	\end{frame}
	
	\begin{frame}{T19: NLS - Initial Guesses}
		\fontsize{6.5}{7.5}\selectfont
		\textbf{TOPIC: Importance of Initial Guesses in NLS}
		
		\vspace{0.1cm}
		\textbf{INITIAL GUESSES:}
		\begin{itemize}
			\item Starting values for parameters
			\item Often small random numbers
			\item Critical for convergence
		\end{itemize}
		
		\vspace{0.1cm}
		\textbf{WHY THEY MATTER:}
		\begin{itemize}
			\item No closed-form solution
			\item Iterative process needs starting point
			\item Different starts may give different results
		\end{itemize}
		
		\vspace{0.1cm}
		\textbf{POTENTIAL ISSUES:}
		\begin{itemize}
			\item Bad starting values → slow convergence
			\item May converge to local minimum
			\item May not converge at all
		\end{itemize}
		
		\vspace{0.1cm}
		\textbf{BEST PRACTICES:}
		\begin{itemize}
			\item Try multiple random starts
			\item Use domain knowledge
			\item Normalize data first
		\end{itemize}
		
		\textcolor{myred}{\textbf{EXAM TRAP:}} NLS starts with \textcolor{myblue}{initial guesses}, not exact values!
	\end{frame}
	
	\begin{frame}{T20: NLS - Convergence Criteria}
		\fontsize{6.5}{7.5}\selectfont
		\textbf{TOPIC: Convergence Criteria in NLS}
		
		\vspace{0.1cm}
		\textbf{DEFINITION:}
		\begin{itemize}
			\item Conditions for stopping the iterative process
			\item Determines when solution is "good enough"
		\end{itemize}
		
		\vspace{0.1cm}
		\textbf{COMMON CRITERIA:}
		\begin{enumerate}
			\item \textcolor{myblue}{Objective improvement:} Change in loss below threshold
			\item \textcolor{myblue}{Gradient magnitude:} Gradient near zero
			\item \textcolor{myblue}{Parameter change:} Parameters stop changing
			\item \textcolor{myblue}{Maximum iterations:} Prevents infinite loops
		\end{enumerate}
		
		\vspace{0.1cm}
		\textbf{TRADE-OFFS:}
		\begin{itemize}
			\item Tight threshold: More iterations, better precision
			\item Loose threshold: Fewer iterations, less precision
		\end{itemize}
		
		\vspace{0.1cm}
		\textbf{EXAMPLE:}
		\begin{itemize}
			\item Stop when gradient < 0.001
			\item Stop after 1000 iterations			\item Stop when loss improves < 0.0001
		\end{itemize}
		
		\textcolor{myred}{\textbf{EXAM TRAP:}} Multiple stopping criteria exist, not just one!
	\end{frame}
	
	% ============================================================
	% SECTION 3: MAXIMUM LIKELIHOOD ESTIMATION
	% ============================================================
	
	\begin{frame}{T21: MLE Definition}
		\fontsize{6.5}{7.5}\selectfont
		\textbf{TOPIC: Maximum Likelihood Estimation - Definition and Overview}
		
		\vspace{0.1cm}
		\textbf{DEFINITION:}
		\begin{itemize}
			\item Probabilistic approach to parameter estimation
			\item Finds parameter values that maximize probability of observing data
			\item More flexible than Least Squares
		\end{itemize}
		
		\vspace{0.1cm}
		\textbf{CORE QUESTION:}
		\begin{itemize}
			\item "Given observed data, what parameter values make this data most probable?"
		\end{itemize}
		
		\vspace{0.1cm}
		\textbf{PROCESS:}
		\begin{enumerate}
			\item Assume data from specific distribution
			\item Construct likelihood function: $L(\theta) = P(X|\theta)$
			\item Find $\theta$ that maximizes $L(\theta)$
			\item Use log-likelihood for easier computation
		\end{enumerate}
		
		\vspace{0.1cm}
		\textbf{KEY PROPERTIES:}
		\begin{itemize}
			\item Consistent
			\item Asymptotically efficient
			\item Asymptotically normal
		\end{itemize}
		
		\textcolor{myred}{\textbf{EXAM TRAP:}} MLE maximizes \textcolor{myblue}{likelihood} = $P(\text{data}|\theta)$!
	\end{frame}
	
	\begin{frame}{T22: MLE vs OLS - Equivalence}
		\fontsize{6.5}{7.5}\selectfont
		\textbf{TOPIC: When OLS and MLE Give Identical Estimates}
		
		\vspace{0.1cm}
		\textbf{EQUIVALENCE CONDITION:}
		\begin{itemize}
			\item Errors are independently normally distributed
			\item Zero mean
			\item Constant variance
		\end{itemize}
		
		\vspace{0.1cm}
		\textbf{WHY EQUIVALENCE:}
		\begin{itemize}
			\item Log-likelihood: $\ell(\beta) = -\frac{N}{2}\log(2\pi\sigma^2) - \frac{1}{2\sigma^2}\sum(y_i - X_i\beta)^2$
			\item Maximizing $\ell(\beta)$ $\equiv$ minimizing $\sum(y_i - X_i\beta)^2$
			\item Same objective function
		\end{itemize}
		
		\vspace{0.1cm}
		\textbf{WHEN THEY DIFFER:}
		\begin{itemize}
			\item Non-normal errors
			\item Heteroscedasticity
			\item Different distributions (Bernoulli, Poisson, etc.)
		\end{itemize}
		
		\vspace{0.1cm}
		\textbf{PRACTICAL IMPLICATION:}
		\begin{itemize}
			\item Can use either method for linear regression
			\item MLE more flexible for other models
		\end{itemize}
		
		\textcolor{myred}{\textbf{EXAM TRAP:}} OLS = MLE under \textcolor{myblue}{normal errors}!
	\end{frame}
	
	\begin{frame}{T23: MLE - Likelihood Function}
		\fontsize{6.5}{7.5}\selectfont
		\textbf{TOPIC: Understanding the Likelihood Function}
		
		\vspace{0.1cm}
		\textbf{DEFINITION:}
		\begin{itemize}
			\item $L(\theta | X) = P(X | \theta)$
			\item Probability of observing data given parameters
			\item Product of individual probabilities (if independent)
		\end{itemize}
		
		\vspace{0.1cm}
		\textbf{FORMULA:}
		\[
		L(\theta) = \prod_{i=1}^N f(x_i | \theta)
		\]
		where $f(x_i|\theta)$ is the probability density/mass function
		
		\vspace{0.1cm}
		\textbf{KEY DISTINCTION:}
		\begin{itemize}
			\item \textcolor{myblue}{Likelihood:} $P(\text{data} | \theta)$
			\item \textcolor{myblue}{Posterior:} $P(\theta | \text{data})$ (Bayesian)
		\end{itemize}
		
		\vspace{0.1cm}
		\textbf{INDEPENDENCE ASSUMPTION:}
		\begin{itemize}
			\item Joint probability = product of individual probabilities
			\item $\prod$ notation for multiplication
		\end{itemize}
		
		\textcolor{myred}{\textbf{EXAM TRAP:}} Likelihood = \textcolor{myblue}{$P(\text{data}|\theta)$}, NOT $P(\theta|\text{data})$!
	\end{frame}
	
	\begin{frame}{T24: MLE - Log-Likelihood}
		\fontsize{6.5}{7.5}\selectfont
		\textbf{TOPIC: Why Use Log-Likelihood Instead of Likelihood}
		
		\vspace{0.1cm}
		\textbf{DEFINITION:}
		\begin{itemize}
			\item Natural logarithm of the likelihood function
			\item $\ell(\theta) = \log L(\theta) = \sum_{i=1}^N \log f(x_i|\theta)$
		\end{itemize}
		
		\vspace{0.1cm}
		\textbf{MATHEMATICAL CONVENIENCE:}
		\begin{itemize}
			\item Likelihood: $L(\theta) = \prod f(x_i|\theta)$
			\item Log-likelihood: $\ell(\theta) = \sum \log f(x_i|\theta)$
			\item Products become sums
		\end{itemize}
		
		\vspace{0.1cm}
		\textbf{MONOTONICITY:}
		\begin{itemize}
			\item $\log$ is strictly increasing
			\item Maximizing $\ell(\theta)$ $\equiv$ maximizing $L(\theta)$
		\end{itemize}
		
		\vspace{0.1cm}
		\textbf{NUMERICAL STABILITY:}
		\begin{itemize}
			\item Products of small probabilities underflow
			\item Sums of logs are stable
		\end{itemize}
		
		\textcolor{myred}{\textbf{EXAM TRAP:}} Log transforms \textcolor{myblue}{products to sums}!
	\end{frame}
	
	\begin{frame}{T25: MLE - Logistic Regression}
		\fontsize{6.5}{7.5}\selectfont
		\textbf{TOPIC: MLE for Logistic Regression}
		
		\vspace{0.1cm}
		\textbf{BINARY OUTCOME:}
		\begin{itemize}
			\item $y_i \in \{0, 1\}$
			\item Bernoulli distribution
		\end{itemize}
		
		\vspace{0.1cm}
		\textbf{BERNOULLI LIKELIHOOD:}
		\[
		L(\beta) = \prod_{i=1}^N p_i^{y_i}(1-p_i)^{1-y_i}
		\]
		where $p_i = \frac{1}{1 + e^{-X_i\beta}}$
		
		\vspace{0.1cm}
		\textbf{LOG-LIKELIHOOD:}
		\[
		\ell(\beta) = \sum_{i=1}^N [y_i \log(p_i) + (1-y_i)\log(1-p_i)]
		\]
		
		\vspace{0.1cm}
		\textbf{MAXIMIZATION:}
		\begin{itemize}
			\item No closed-form solution
			\item Numerical optimization (Newton-Raphson, IRLS)
		\end{itemize}
		
		\textcolor{myred}{\textbf{EXAM TRAP:}} Logistic regression uses \textcolor{myblue}{Bernoulli} distribution!
	\end{frame}
	
	\begin{frame}{T26: MLE - Linear Regression}
		\fontsize{6.5}{7.5}\selectfont
		\textbf{TOPIC: MLE for Linear Regression}
		
		\vspace{0.1cm}
		\textbf{CONTINUOUS OUTCOME:}
		\begin{itemize}
			\item $y_i \in \mathbb{R}$ (real numbers)
			\item Normal distribution for errors
		\end{itemize}
		
		\vspace{0.1cm}
		\textbf{NORMAL LIKELIHOOD:}
		\[
		L(\beta, \sigma^2) = \prod_{i=1}^N \frac{1}{\sqrt{2\pi\sigma^2}} \exp\left(-\frac{(y_i - X_i\beta)^2}{2\sigma^2}\right)
		\]
		
		\vspace{0.1cm}
		\textbf{LOG-LIKELIHOOD:}
		\[
		\ell(\beta) = -\frac{N}{2}\log(2\pi\sigma^2) - \frac{1}{2\sigma^2}\sum(y_i - X_i\beta)^2
		\]
		
		\vspace{0.1cm}
		\textbf{EQUIVALENCE TO OLS:}
		\begin{itemize}
			\item Maximizing $\ell$ = minimizing RSS
			\item Same parameter estimates
		\end{itemize}
		
		\textcolor{myred}{\textbf{EXAM TRAP:}} Linear regression uses \textcolor{myblue}{Normal} distribution!
	\end{frame}
	
	\begin{frame}{T27: MLE - Logit Model}
		\fontsize{6.5}{7.5}\selectfont
		\textbf{TOPIC: Log-Likelihood Function for Logit Model}
		
		\vspace{0.1cm}
		\textbf{LOGIT MODEL:}
		\begin{itemize}
			\item $P_i = \frac{1}{1 + e^{-(\beta_0 + \beta_1 x_{1i} + ... + \beta_m x_{mi})}}$
			\item $y_i = 1$ with probability $P_i$
			\item $y_i = 0$ with probability $1-P_i$
		\end{itemize}
		
		\vspace{0.1cm}
		\textbf{LIKELIHOOD FOR ONE OBSERVATION:}
		\[
		L(y_i) = F(y_i)^{y_i}(1-F(y_i))^{1-y_i}
		\]
		
		\vspace{0.1cm}
		\textbf{JOINT LIKELIHOOD:}
		\[
		L = \prod_{i=1}^N F(y_i)^{y_i}(1-F(y_i))^{1-y_i}
		\]
		
		\vspace{0.1cm}
		\textbf{LOG-LIKELIHOOD:}
		\[
		\log L = \sum_{i=1}^N [y_i \log(F(y_i)) + (1-y_i)\log(1-F(y_i))]
		\]
		
		\textcolor{myred}{\textbf{EXAM TRAP:}} Logit log-likelihood uses \textcolor{myblue}{log of F and 1-F}!
	\end{frame}
	
	\begin{frame}{T28: MLE - Threshold for Classification}
		\fontsize{6.5}{7.5}\selectfont
		\textbf{TOPIC: Threshold Z in Classification}
		
		\vspace{0.1cm}
		\textbf{DEFINITION:}
		\begin{itemize}
			\item Threshold Z converts predicted probabilities into class predictions
			\item Classification rule: $\hat{y}_i = 1$ if $P_i \geq Z$, else $\hat{y}_i = 0$
		\end{itemize}
		
		\vspace{0.1cm}
		\textbf{COMMON THRESHOLDS:}
		\begin{itemize}
			\item \textcolor{myblue}{Z = 0.5:} Balanced costs
			\item \textcolor{myblue}{Lower Z:} More conservative (predict more positives)
			\item \textcolor{myblue}{Higher Z:} More liberal (predict fewer positives)
		\end{itemize}
		
		\vspace{0.1cm}
		\textbf{EXAMPLE (LOAN DEFAULT):}
		\begin{itemize}
			\item Z = 0.05 (5\%)
			\item Because cost of bad loan > profit foregone
			\item More conservative lending decisions
		\end{itemize}
		
		\vspace{0.1cm}
		\textbf{KEY INSIGHT:}
		\begin{itemize}
			\item Threshold depends on relative costs
			\item Lower threshold when false negatives more costly
		\end{itemize}
		
		\textcolor{myred}{\textbf{EXAM TRAP:}} Threshold converts \textcolor{myblue}{probabilities to class labels}!
	\end{frame}
	
	\begin{frame}{T29: MLE - Cost Asymmetry}
		\fontsize{6.5}{7.5}\selectfont
		\textbf{TOPIC: Cost Asymmetry and Threshold Selection}
		
		\vspace{0.1cm}
		\textbf{TYPE I ERROR (FALSE POSITIVE):}
		\begin{itemize}
			\item Predict default, actually no default
			\item Cost = Lost profit from not making loan
		\end{itemize}
		
		\vspace{0.1cm}
		\textbf{TYPE II ERROR (FALSE NEGATIVE):}
		\begin{itemize}
			\item Predict no default, actually defaults
			\item Cost = Actual loan loss (much larger)
		\end{itemize}
		
		\vspace{0.1cm}
		\textbf{RATIONALE FOR LOW THRESHOLD:}
		\begin{itemize}
			\item More conservative (predict more defaults)
			\item Catches more actual defaults
			\item Accepts more false positives
		\end{itemize}
		
		\vspace{0.1cm}
		\textbf{GENERAL PRINCIPLE:}
		\begin{itemize}
			\item Threshold depends on relative costs
			\item Lower threshold when FN more costly
		\end{itemize}
		
		\textcolor{myred}{\textbf{EXAM TRAP:}} Lower threshold when \textcolor{myblue}{missing positive is costly}!
	\end{frame}
	
	\begin{frame}{T30: MLE vs OLS - Comparison}
		\fontsize{6.5}{7.5}\selectfont
		\textbf{TOPIC: Key Advantages of MLE Over OLS}
		
		\vspace{0.1cm}
		\textbf{MLE FLEXIBILITY:}
		\begin{itemize}
			\item Works with any distribution
			\item Bernoulli (logistic regression)
			\item Normal (linear regression)
			\item Poisson (count data)
			\item GARCH (volatility)
		\end{itemize}
		
		\vspace{0.1cm}
		\textbf{OLS LIMITATIONS:}
		\begin{itemize}
			\item Only for continuous outcomes
			\item Assumes normal errors
			\item Cannot handle binary outcomes
		\end{itemize}
		
		\vspace{0.1cm}
		\textbf{WHEN TO USE EACH:}
		\begin{itemize}
			\item \textcolor{myblue}{OLS:} Simple linear regression
			\item \textcolor{myblue}{MLE:} Complex probabilistic models
		\end{itemize}
		
		\textcolor{myred}{\textbf{EXAM TRAP:}} MLE is \textcolor{myblue}{more flexible} than OLS!
	\end{frame}
	
	\begin{frame}{T31: MLE - Binary Outcomes}
		\fontsize{6.5}{7.5}\selectfont
		\textbf{TOPIC: Why OLS Cannot Handle Binary Outcomes}
		
		\vspace{0.1cm}
		\textbf{BINARY OUTCOME ISSUES:}
		\begin{itemize}
			\item $y_i \in \{0, 1\}$ (not continuous)
			\item Errors not normally distributed
			\item Predictions can be outside [0, 1]
		\end{itemize}
		
		\vspace{0.1cm}
		\textbf{PROBLEMS WITH OLS:}
		\begin{itemize}
			\item Heteroscedastic errors
			\item Inefficient estimates
			\item Invalid predictions
		\end{itemize}
		
		\vspace{0.1cm}
		\textbf{SOLUTION:}
		\begin{itemize}
			\item Use logistic regression (MLE)
			\item Predict probabilities bounded [0, 1]
		\end{itemize}
		
		\vspace{0.1cm}
		\textbf{EXAMPLE:}
		\begin{itemize}
			\item Loan default: 0 = no default, 1 = default
			\item OLS would predict values like -0.3 or 1.5
			\item Logistic regression predicts probabilities
		\end{itemize}
		
		\textcolor{myred}{\textbf{EXAM TRAP:}} OLS inappropriate for \textcolor{myblue}{binary outcomes}!
	\end{frame}
	
	\begin{frame}{T32: MLE - Log-Likelihood Formula}
		\fontsize{6.5}{7.5}\selectfont
		\textbf{TOPIC: Log-Likelihood Formula for Logistic Regression}
		
		\vspace{0.1cm}
		\textbf{FORMULA:}
		\[
		\ell = \sum_{i=1}^N [y_i \log(p_i) + (1-y_i)\log(1-p_i)]
		\]
		where $p_i = \frac{1}{1 + e^{-X_i\beta}}$
		
		\vspace{0.1cm}
		\textbf{INTERPRETATION:}
		\begin{itemize}
			\item For $y_i = 1$: $\log(p_i)$
			\item For $y_i = 0$: $\log(1-p_i)$
		\end{itemize}
		
		\vspace{0.1cm}
		\textbf{MAXIMIZATION:}
		\begin{itemize}
			\item Numerical optimization needed
			\item No closed-form solution
			\item Newton-Raphson or IRLS
		\end{itemize}
		
		\vspace{0.1cm}
		\textbf{PROPERTIES:}
		\begin{itemize}
			\item Concave function
			\item Unique maximum
			\item Fast convergence
		\end{itemize}
		
		\textcolor{myred}{\textbf{EXAM TRAP:}} Know the \textcolor{myblue}{log-likelihood formula} for logistic regression!
	\end{frame}
	
	\begin{frame}{T33: MLE - Optimization}
		\fontsize{6.5}{7.5}\selectfont
		\textbf{TOPIC: Optimization Methods for MLE in Logistic Regression}
		
		\vspace{0.1cm}
		\textbf{NO CLOSED-FORM:}
		\begin{itemize}
			\item Log-likelihood is nonlinear in parameters
			\item Cannot solve analytically
		\end{itemize}
		
		\vspace{0.1cm}
		\textbf{NUMERICAL METHODS:}
		\begin{itemize}
			\item \textcolor{myblue}{Newton-Raphson:} Uses second derivatives (Hessian)
			\item \textcolor{myblue}{IRLS:} Iteratively Reweighted Least Squares
			\item \textcolor{myblue}{Gradient Descent:} First-order method
		\end{itemize}
		
		\vspace{0.1cm}
		\textbf{CONVERGENCE:}
		\begin{itemize}
			\item Usually converges quickly
			\item May fail with perfect separation
		\end{itemize}
		
		\vspace{0.1cm}
		\textbf{PERFECT SEPARATION:}
		\begin{itemize}
			\item When classes perfectly separated
			\item Coefficients diverge to infinity
			\item Need regularization
		\end{itemize}
		
		\textcolor{myred}{\textbf{EXAM TRAP:}} Logistic regression requires \textcolor{myblue}{numerical optimization}!
	\end{frame}
	
	\begin{frame}{T34: MLE - Consistency}
		\fontsize{6.5}{7.5}\selectfont
		\textbf{TOPIC: Consistency Property of MLE}
		
		\vspace{0.1cm}
		\textbf{DEFINITION:}
		\begin{itemize}
			\item Estimator converges to true parameter as sample size increases
			\item $\hat{\theta}_{MLE} \xrightarrow{p} \theta_0$ as $n \to \infty$
		\end{itemize}
		
		\vspace{0.1cm}
		\textbf{MLE PROPERTIES:}
		\begin{itemize}
			\item \textcolor{myblue}{Consistent:} Converges to true value
			\item \textcolor{myblue}{Asymptotically efficient:} Minimum variance
			\item \textcolor{myblue}{Asymptotically normal:} Normal distribution
		\end{itemize}
		
		\vspace{0.1cm}
		\textbf{CAVEAT:}
		\begin{itemize}
			\item MLE can be biased in finite samples
			\item Example: $\hat{\sigma}^2_{MLE}$ is biased for $\sigma^2$
			\item Consistency is asymptotic property
		\end{itemize}
		
		\vspace{0.1cm}
		\textbf{WHY IT MATTERS:}
		\begin{itemize}
			\item Reliable estimates with large samples
			\item Foundation for inference
		\end{itemize}
		
		\textcolor{myred}{\textbf{EXAM TRAP:}} MLE is \textcolor{myblue}{consistent} but NOT necessarily unbiased!
	\end{frame}
	
	\begin{frame}{T35: MLE - Asymptotic Normality}
		\fontsize{6.5}{7.5}\selectfont
		\textbf{TOPIC: Asymptotic Distribution of MLE}
		
		\vspace{0.1cm}
		\textbf{FORMULA:}
		\[
		\sqrt{n}(\hat{\theta}_{MLE} - \theta_0) \xrightarrow{d} N(0, I^{-1}(\theta_0))
		\]
		where $I(\theta)$ = Fisher Information matrix
		
		\vspace{0.1cm}
		\textbf{IMPLICATIONS:}
		\begin{itemize}
			\item Confidence intervals can be constructed
			\item Hypothesis tests can be performed
			\item MLE achieves efficiency bound
		\end{itemize}
		
		\vspace{0.1cm}
		\textbf{FISHER INFORMATION:}
		\[
		I(\theta) = -E\left[\frac{\partial^2 \ell(\theta)}{\partial \theta^2}\right]
		\]
		
		\vspace{0.1cm}
		\textbf{PRACTICAL USE:}
		\begin{itemize}
			\item Standard errors from inverse Hessian
			\item Wald tests, likelihood ratio tests
		\end{itemize}
		
		\textcolor{myred}{\textbf{EXAM TRAP:}} MLE is \textcolor{myblue}{asymptotically normal}!
	\end{frame}
	
	\begin{frame}{T36: MLE - Efficiency}
		\fontsize{6.5}{7.5}\selectfont
		\textbf{TOPIC: Asymptotic Efficiency of MLE}
		
		\vspace{0.1cm}
		\textbf{DEFINITION:}
		\begin{itemize}
			\item MLE achieves lowest possible variance among consistent estimators asymptotically
			\item No consistent estimator has lower variance
		\end{itemize}
		
		\vspace{0.1cm}
		\textbf{CRAMÉR-RAO LOWER BOUND:}
		\begin{itemize}
			\item Minimum variance for unbiased estimators
			\item MLE achieves this asymptotically
			\item $Var(\hat{\theta}) \geq \frac{1}{I(\theta)}$
		\end{itemize}
		
		\vspace{0.1cm}
		\textbf{CAVEATS:}
		\begin{itemize}
			\item Finite sample properties may differ
			\item Regularity conditions must hold
		\end{itemize}
		
		\vspace{0.1cm}
		\textbf{WHY IT MATTERS:}
		\begin{itemize}
			\item Most precise estimates possible
			\item Efficient use of data
		\end{itemize}
		
		\textcolor{myred}{\textbf{EXAM TRAP:}} MLE is \textcolor{myblue}{asymptotically efficient}!
	\end{frame}
	
	\begin{frame}{T37: MLE - Invariance Property}
		\fontsize{6.5}{7.5}\selectfont
		\textbf{TOPIC: Invariance Property of MLE}
		
		\vspace{0.1cm}
		\textbf{DEFINITION:}
		\begin{itemize}
			\item If $\hat{\theta}_{MLE}$ is MLE of $\theta$
			\item Then $g(\hat{\theta}_{MLE})$ is MLE of $g(\theta)$
			\item For any continuous function $g$
		\end{itemize}
		
		\vspace{0.1cm}
		\textbf{EXAMPLE:}
		\begin{itemize}
			\item If $\hat{\sigma}^2$ is MLE of $\sigma^2$
			\item Then $\sqrt{\hat{\sigma}^2}$ is MLE of $\sigma$
		\end{itemize}
		
		\vspace{0.1cm}
		\textbf{WHY USEFUL:}
		\begin{itemize}
			\item Easy transformation of estimates
			\item No need to re-estimate for transformed parameters
		\end{itemize}
		
		\vspace{0.1cm}
		\textbf{IMPORTANT:}
		\begin{itemize}
			\item Invariance does NOT hold for expected values
			\item $E[g(\hat{\theta})] \neq g(E[\hat{\theta}])$ generally
		\end{itemize}
		
		\textcolor{myred}{\textbf{EXAM TRAP:}} MLE is \textcolor{myblue}{invariant} to parameter transformations!
	\end{frame}
	
	\begin{frame}{T38: MLE - Fisher Information}
		\fontsize{6.5}{7.5}\selectfont
		\textbf{TOPIC: Fisher Information in MLE}
		
		\vspace{0.1cm}
		\textbf{DEFINITION:}
		\begin{itemize}
			\item Measure of amount of information data carries about unknown parameter
			\item Higher information = more precise estimates
		\end{itemize}
		
		\vspace{0.1cm}
		\textbf{FORMULA:}
		\[
		I(\theta) = -E\left[\frac{\partial^2 \ell(\theta)}{\partial \theta^2}\right]
		\]
		
		\vspace{0.1cm}
		\textbf{CONNECTION TO VARIANCE:}
		\begin{itemize}
			\item $Var(\hat{\theta}) \geq \frac{1}{I(\theta)}$
			\item Cramér-Rao lower bound
		\end{itemize}
		
		\vspace{0.1cm}
		\textbf{INTERPRETATION:}
		\begin{itemize}
			\item More data = more information
			\item Sharper likelihood peak = more information
		\end{itemize}
		
		\textcolor{myred}{\textbf{EXAM TRAP:}} Fisher Information measures \textcolor{myblue}{information about parameters}!
	\end{frame}
	
	\begin{frame}{T39: MLE - Applications}
		\fontsize{6.5}{7.5}\selectfont
		\textbf{TOPIC: Applications of MLE}
		
		\vspace{0.1cm}
		\textbf{LINEAR REGRESSION:}
		\begin{itemize}
			\item Can use OLS or MLE
			\item Equivalent under normal errors
		\end{itemize}
		
		\vspace{0.1cm}
		\textbf{LOGISTIC REGRESSION:}
		\begin{itemize}
			\item Uses MLE (Bernoulli likelihood)
			\item No OLS equivalent
		\end{itemize}
		
		\vspace{0.1cm}
		\textbf{OTHER MLE APPLICATIONS:}
		\begin{itemize}
			\item Probit models
			\item GARCH models
			\item Count data models
			\item Survival analysis
		\end{itemize}
		
		\vspace{0.1cm}
		\textbf{KEY INSIGHT:}
		\begin{itemize}
			\item MLE is flexible framework
			\item Works with any distribution
			\item Foundation for many models
		\end{itemize}
		
		\textcolor{myred}{\textbf{EXAM TRAP:}} Both linear and logistic regression can use \textcolor{myblue}{MLE}!
	\end{frame}
	
	\begin{frame}{T40: MLE - Summary}
		\fontsize{6.5}{7.5}\selectfont
		\textbf{TOPIC: Key Advantages of MLE Over OLS}
		
		\vspace{0.1cm}
		\textbf{FLEXIBILITY OF MLE:}
		\begin{itemize}
			\item Works with any distribution
			\item Handles binary, count, continuous outcomes
			\item Used in complex models (GARCH, neural networks)
		\end{itemize}
		
		\vspace{0.1cm}
		\textbf{LIMITATIONS OF OLS:}
		\begin{itemize}
			\item Only for continuous outcomes
			\item Assumes normal errors
			\item Cannot handle binary outcomes
		\end{itemize}
		
		\vspace{0.1cm}
		\textbf{PRACTICAL GUIDANCE:}
		\begin{itemize}
			\item \textcolor{myblue}{OLS:} Simple linear regression
			\item \textcolor{myblue}{MLE:} Complex probabilistic models
		\end{itemize}
		
		\vspace{0.1cm}
		\textbf{REMEMBER:}
		\begin{itemize}
			\item MLE more flexible but more complex
			\item OLS simpler for basic regression
		\end{itemize}
		
		\textcolor{myred}{\textbf{EXAM TRAP:}} MLE is \textcolor{myblue}{more flexible} than OLS!
	\end{frame}
	
	% ============================================================
	% SECTION 4: GRADIENT DESCENT
	% ============================================================
	
	\begin{frame}{T41: Gradient Descent Definition}
		\fontsize{6.5}{7.5}\selectfont
		\textbf{TOPIC: Gradient Descent - Definition and Overview}
		
		\vspace{0.1cm}
		\textbf{DEFINITION:}
		\begin{itemize}
			\item Iterative optimization algorithm
			\item Updates parameters in direction of negative gradient
			\item Workhorse of modern machine learning
		\end{itemize}
		
		\vspace{0.1cm}
		\textbf{CORE CONCEPT:}
		\begin{itemize}
			\item Move in direction of steepest descent
			\item Update rule: $w^{new} = w^{old} - \eta \frac{\partial L}{\partial w}$
			\item $\eta$ = learning rate
			\item $\frac{\partial L}{\partial w}$ = gradient
		\end{itemize}
		
		\vspace{0.1cm}
		\textbf{PROCESS:}
		\begin{enumerate}
			\item Start with initial guesses
			\item Calculate gradient
			\item Update parameters in opposite direction
			\item Repeat until convergence
		\end{enumerate}
		
		\vspace{0.1cm}
		\textbf{CONVERGENCE:}
		\begin{itemize}
			\item Gradient near zero
			\item Loss no longer decreasing
		\end{itemize}
		
		\textcolor{myred}{\textbf{EXAM TRAP:}} Gradient descent uses \textcolor{myblue}{negative gradient} direction!
	\end{frame}
	
	\begin{frame}{T42: Gradient Descent - Update Rule}
		\fontsize{6.5}{7.5}\selectfont
		\textbf{TOPIC: Gradient Descent Update Rule}
		
		\vspace{0.1cm}
		\textbf{FORMULA:}
		\[
		w_i^{new} = w_i^{old} - \eta \frac{\partial \mathcal{L}(w_i)}{\partial w_i}
		\]
		
		\vspace{0.1cm}
		\textbf{COMPONENTS:}
		\begin{itemize}
			\item $w_i$: Parameter to update
			\item $\eta$: Learning rate (0, 1)
			\item $\frac{\partial \mathcal{L}}{\partial w_i}$: Partial derivative (gradient)
		\end{itemize}
		
		\vspace{0.1cm}
		\textbf{WHY SUBTRACT:}
		\begin{itemize}
			\item Gradient points uphill
			\item Negative gradient points downhill
			\item Subtracting moves toward minimum
		\end{itemize}
		
		\vspace{0.1cm}
		\textbf{EXAMPLE:}
		\begin{itemize}
			\item Current $w = 0.8$, gradient = -48.66, $\eta = 0.005$
			\item Update: $0.8 - 0.005(-48.66) = 1.043$
		\end{itemize}
		
		\textcolor{myred}{\textbf{EXAM TRAP:}} Gradient descent \textcolor{myblue}{subtracts} gradient!
	\end{frame}
	
	\begin{frame}{T43: Gradient Descent - Learning Rate}
		\fontsize{6.5}{7.5}\selectfont
		\textbf{TOPIC: Learning Rate - Too Small}
		
		\vspace{0.1cm}
		\textbf{SMALL LEARNING RATE:}
		\begin{itemize}
			\item Small steps
			\item Slow convergence
			\item Many iterations needed
			\item May get stuck in local minima
		\end{itemize}
		
		\vspace{0.1cm}
		\textbf{VISUAL:}
		\begin{itemize}
			\item Tiny steps down the valley
			\item Takes forever to reach bottom
		\end{itemize}
		
		\vspace{0.1cm}
		\textbf{SOLUTION:}
		\begin{itemize}
			\item Increase learning rate
			\item Use dynamic learning rate
		\end{itemize}
		
		\vspace{0.1cm}
		\textbf{EXAMPLE:}
		\begin{itemize}
			\item $\eta = 0.0001$: Very slow
			\item $\eta = 0.001$: Still slow
			\item $\eta = 0.01$: Reasonable
		\end{itemize}
		
		\textcolor{myred}{\textbf{EXAM TRAP:}} Too small $\eta$ = \textcolor{myblue}{slow convergence}!
	\end{frame}
	
	\begin{frame}{T44: Gradient Descent - Too Large Learning Rate}
		\fontsize{6.5}{7.5}\selectfont
		\textbf{TOPIC: Learning Rate - Too Large}
		
		\vspace{0.1cm}
		\textbf{LARGE LEARNING RATE:}
		\begin{itemize}
			\item Large steps
			\item Overshooting minimum
			\item Oscillation around minimum
			\item Divergence (loss increases)
		\end{itemize}
		
		\vspace{0.1cm}
		\textbf{VISUAL:}
		\begin{itemize}
			\item Jumping over the valley
			\item Bouncing back and forth
		\end{itemize}
		
		\vspace{0.1cm}
		\textbf{SOLUTION:}
		\begin{itemize}
			\item Decrease learning rate
			\item Use dynamic learning rate
		\end{itemize}
		
		\vspace{0.1cm}
		\textbf{EXAMPLE:}
		\begin{itemize}
			\item $\eta = 10$: Diverges
			\item $\eta = 1$: Oscillates
			\item $\eta = 0.1$: Maybe OK
		\end{itemize}
		
		\textcolor{myred}{\textbf{EXAM TRAP:}} Too large $\eta$ = \textcolor{myblue}{overshooting/divergence}!
	\end{frame}
	
	\begin{frame}{T45: Gradient Descent - Epoch}
		\fontsize{6.5}{7.5}\selectfont
		\textbf{TOPIC: Epoch in Gradient Descent}
		
		\vspace{0.1cm}
		\textbf{DEFINITION:}
		\begin{itemize}
			\item One complete pass through the entire training dataset
			\item Includes calculating loss and adjusting weights
		\end{itemize}
		
		\vspace{0.1cm}
		\textbf{BATCH GRADIENT DESCENT:}
		\begin{itemize}
			\item One epoch = one weight update
			\item Uses all data for each update
		\end{itemize}
		
		\vspace{0.1cm}
		\textbf{STOCHASTIC GRADIENT DESCENT:}
		\begin{itemize}
			\item One epoch = N weight updates
			\item Uses one data point per update
		\end{itemize}
		
		\vspace{0.1cm}
		\textbf{EXAMPLE:}
		\begin{itemize}
			\item Training data: 1000 observations
			\item Batch GD: 1 update per epoch
			\item SGD: 1000 updates per epoch
		\end{itemize}
		
		\textcolor{myred}{\textbf{EXAM TRAP:}} Epoch = \textcolor{myblue}{one complete pass} through data!
	\end{frame}
	
	\begin{frame}{T46: Gradient Descent - Batch}
		\fontsize{6.5}{7.5}\selectfont
		\textbf{TOPIC: Batch Gradient Descent}
		
		\vspace{0.1cm}
		\textbf{DEFINITION:}
		\begin{itemize}
			\item Uses entire training dataset for each update
			\item Computes gradient over all N samples
		\end{itemize}
		
		\vspace{0.1cm}
		\textbf{CHARACTERISTICS:}
		\begin{itemize}
			\item One update per epoch
			\item Stable convergence
			\item Deterministic
		\end{itemize}
		
		\vspace{0.1cm}
		\textbf{ADVANTAGES:}
		\begin{itemize}
			\item Stable convergence
			\item Deterministic
		\end{itemize}
		
		\vspace{0.1cm}
		\textbf{DISADVANTAGES:}
		\begin{itemize}
			\item Computationally expensive
			\item Requires all data in memory
		\end{itemize}
		
		\textcolor{myred}{\textbf{EXAM TRAP:}} Batch GD uses \textcolor{myblue}{entire dataset} per update!
	\end{frame}
	
	\begin{frame}{T47: Gradient Descent - Stochastic}
		\fontsize{6.5}{7.5}\selectfont
		\textbf{TOPIC: Stochastic Gradient Descent (SGD)}
		
		\vspace{0.1cm}
		\textbf{DEFINITION:}
		\begin{itemize}
			\item Uses one randomly selected data point per update
			\item Updates weights after each observation
		\end{itemize}
		
		\vspace{0.1cm}
		\textbf{CHARACTERISTICS:}
		\begin{itemize}
			\item N updates per epoch
			\item Noisy convergence
			\item Random selection
		\end{itemize}
		
		\vspace{0.1cm}
		\textbf{ADVANTAGES:}
		\begin{itemize}
			\item Lower memory requirements
			\item Faster per iteration
			\item Can escape local minima (noise)
		\end{itemize}
		
		\vspace{0.1cm}
		\textbf{DISADVANTAGES:}
		\begin{itemize}
			\item Noisy convergence
			\item More iterations needed
		\end{itemize}
		
		\textcolor{myred}{\textbf{EXAM TRAP:}} SGD uses \textcolor{myblue}{one data point} per update!
	\end{frame}
	
	\begin{frame}{T48: Gradient Descent - Mini-Batch}
		\fontsize{6.5}{7.5}\selectfont
		\textbf{TOPIC: Mini-Batch Gradient Descent}
		
		\vspace{0.1cm}
		\textbf{DEFINITION:}
		\begin{itemize}
			\item Uses small subset (batch) of data for each update
			\item Balances batch and stochastic methods
		\end{itemize}
		
		\vspace{0.1cm}
		\textbf{CHARACTERISTICS:}
		\begin{itemize}
			\item Batch size: 32, 64, 128
			\item Multiple updates per epoch
			\item Most common in practice
		\end{itemize}
		
		\vspace{0.1cm}
		\textbf{ADVANTAGES:}
		\begin{itemize}
			\item More stable than SGD
			\item Faster than batch
			\item Efficient memory usage
		\end{itemize}
		
		\vspace{0.1cm}
		\textbf{COMMON IN PRACTICE:}
		\begin{itemize}
			\item Standard for deep learning
			\item Used in most frameworks
		\end{itemize}
		
		\textcolor{myred}{\textbf{EXAM TRAP:}} Mini-batch uses \textcolor{myblue}{small subset} of data!
	\end{frame}
	
	\begin{frame}{T49: Gradient Descent - Convergence}
		\fontsize{6.5}{7.5}\selectfont
		\textbf{TOPIC: Convergence in Gradient Descent}
		
		\vspace{0.1cm}
		\textbf{DEFINITION:}
		\begin{itemize}
			\item Algorithm has reached a point where further updates don't improve loss significantly
			\item Gradient is near zero
		\end{itemize}
		
		\vspace{0.1cm}
		\textbf{CONVERGENCE CRITERIA:}
		\begin{enumerate}
			\item Gradient magnitude below threshold
			\item Loss improvement below threshold
			\item Maximum iterations reached
			\item Validation error starts increasing
		\end{enumerate}
		
		\vspace{0.1cm}
		\textbf{PRACTICAL CONSIDERATIONS:}
		\begin{itemize}
			\item Exact zero gradient rare
			\item Tradeoff: precision vs computation
			\item Early stopping prevents overfitting
		\end{itemize}
		
		\vspace{0.1cm}
		\textbf{EXAMPLE:}
		\begin{itemize}
			\item Stop when gradient < 0.001
			\item Stop after 1000 iterations
		\end{itemize}
		
		\textcolor{myred}{\textbf{EXAM TRAP:}} Convergence = \textcolor{myblue}{improvement below threshold}, not exactly zero gradient!
	\end{frame}
	
	\begin{frame}{T50: Gradient Descent - Local Minima}
		\fontsize{6.5}{7.5}\selectfont
		\textbf{TOPIC: Local Minima in Gradient Descent}
		
		\vspace{0.1cm}
		\textbf{DEFINITION:}
		\begin{itemize}
			\item Point where loss is lower than neighbors but not lowest possible
			\item Gradient = 0
			\item Not global minimum
		\end{itemize}
		
		\vspace{0.1cm}
		\textbf{PROBLEM:}
		\begin{itemize}
			\item Gradient descent can get stuck
			\item May not find global minimum
		\end{itemize}
		
		\vspace{0.1cm}
		\textbf{SOLUTIONS:}
		\begin{itemize}
			\item Momentum
			\item Multiple restarts
			\item Stochastic gradient descent
			\item Adaptive learning rates
		\end{itemize}
		
		\vspace{0.1cm}
		\textbf{IN HIGH DIMENSIONS:}
		\begin{itemize}
			\item Saddle points more common than local minima
			\item Neural networks: local minima rare
		\end{itemize}
		
		\textcolor{myred}{\textbf{EXAM TRAP:}} Local minima are \textcolor{myblue}{traps} for gradient descent!
	\end{frame}
	
	\begin{frame}{T51: Gradient Descent - Momentum}
		\fontsize{6.5}{7.5}\selectfont
		\textbf{TOPIC: Momentum in Gradient Descent}
		
		\vspace{0.1cm}
		\textbf{DEFINITION:}
		\begin{itemize}
			\item Accelerates convergence by accumulating past gradients
			\item Reduces oscillations
		\end{itemize}
		
		\vspace{0.1cm}
		\textbf{UPDATE RULE:}
		\[
		w^{new} = w^{old} - \eta \frac{\partial \mathcal{L}}{\partial w^{old}} + \mu(w^{old} - w^{older})
		\]
		where $\mu$ = momentum rate (0, 1)
		
		\vspace{0.1cm}
		\textbf{BENEFITS:}
		\begin{itemize}
			\item Speeds up convergence
			\item Reduces oscillation
			\item Helps escape local minima
		\end{itemize}
		
		\vspace{0.1cm}
		\textbf{INTUITION:}
		\begin{itemize}
			\item Like a ball rolling down hill
			\item Builds up speed in consistent directions
			\item Dampens oscillations
		\end{itemize}
		
		\textcolor{myred}{\textbf{EXAM TRAP:}} Momentum uses \textcolor{myblue}{past gradients} to smooth updates!
	\end{frame}
	
	\begin{frame}{T52: Gradient Descent - Dynamic Learning Rate}
		\fontsize{6.5}{7.5}\selectfont
		\textbf{TOPIC: Dynamic Learning Rate}
		
		\vspace{0.1cm}
		\textbf{DEFINITION:}
		\begin{itemize}
			\item Learning rate changes during training
			\item Start with larger $\eta$, reduce over time
		\end{itemize}
		
		\vspace{0.1cm}
		\textbf{DECAY FUNCTIONS:}
		\begin{itemize}
			\item Exponential: $\eta_t = \eta_0 e^{-kt}$
			\item Inverse: $\eta_t = \frac{\eta_0}{1 + kt}$
		\end{itemize}
		
		\vspace{0.1cm}
		\textbf{BENEFITS:}
		\begin{itemize}
			\item Fast initial progress
			\item Fine-tuning near optimum
			\item Better convergence
		\end{itemize}
		
		\vspace{0.1cm}
		\textbf{EXAMPLE:}
		\begin{itemize}
			\item Start: $\eta_0 = 0.1$
			\item After 100 epochs: $\eta = 0.01$
			\item After 1000 epochs: $\eta = 0.001$
		\end{itemize}
		
		\textcolor{myred}{\textbf{EXAM TRAP:}} Dynamic learning rate \textcolor{myblue}{decays} over time!
	\end{frame}
	
	\begin{frame}{T53: Gradient Descent - Exponential Decay}
		\fontsize{6.5}{7.5}\selectfont
		\textbf{TOPIC: Exponential Decay for Learning Rate}
		
		\vspace{0.1cm}
		\textbf{FORMULA:}
		\[
		\eta_t = \eta_0 e^{-kt}
		\]
		
		\vspace{0.1cm}
		\textbf{COMPONENTS:}
		\begin{itemize}
			\item $\eta_0$: Initial learning rate
			\item $k$: Decay rate
			\item $t$: Epoch number
		\end{itemize}
		
		\vspace{0.1cm}
		\textbf{PROPERTIES:}
		\begin{itemize}
			\item Smooth continuous decay
			\item Never reaches exactly zero
			\item Common in practice
		\end{itemize}
		
		\vspace{0.1cm}
		\textbf{EXAMPLE:}
		\begin{itemize}
			\item $\eta_0 = 0.1$, $k = 0.01$
			\item $t = 0$: $\eta = 0.1$
			\item $t = 100$: $\eta = 0.1 e^{-1} = 0.037$
			\item $t = 200$: $\eta = 0.1 e^{-2} = 0.014$
		\end{itemize}
		
		\textcolor{myred}{\textbf{EXAM TRAP:}} Know the \textcolor{myblue}{exponential decay formula}!
	\end{frame}
	
	\begin{frame}{T54: Gradient Descent - Inverse Decay}
		\fontsize{6.5}{7.5}\selectfont
		\textbf{TOPIC: Inverse Decay for Learning Rate}
		
		\vspace{0.1cm}
		\textbf{FORMULA:}
		\[
		\eta_t = \frac{\eta_0}{1 + kt}
		\]
		
		\vspace{0.1cm}
		\textbf{COMPONENTS:}
		\begin{itemize}
			\item $\eta_0$: Initial learning rate
			\item $k$: Decay rate
			\item $t$: Epoch number
		\end{itemize}
		
		\vspace{0.1cm}
		\textbf{PROPERTIES:}
		\begin{itemize}
			\item Slower decay than exponential
			\item Approaches zero asymptotically
		\end{itemize}
		
		\vspace{0.1cm}
		\textbf{EXAMPLE:}
		\begin{itemize}
			\item $\eta_0 = 0.1$, $k = 0.01$
			\item $t = 0$: $\eta = 0.1$
			\item $t = 100$: $\eta = 0.1/2 = 0.05$
			\item $t = 200$: $\eta = 0.1/3 = 0.033$
		\end{itemize}
		
		\textcolor{myred}{\textbf{EXAM TRAP:}} Inverse decay: $\eta_t = \frac{\eta_0}{1 + kt}$!
	\end{frame}
	
	\begin{frame}{T55: Gradient Descent - Example}
		\fontsize{6.5}{7.5}\selectfont
		\textbf{TOPIC: Gradient Descent Example - Salary Data}
		
		\vspace{0.1cm}
		\textbf{EXAMPLE SETUP:}
		\begin{itemize}
			\item Known: $b_0 = 16.88$
			\item Optimize: $b_1$
			\item Initial: $b_1 = 0.8$
			\item Learning rate: $\eta = 0.005$
		\end{itemize}
		
		\vspace{0.1cm}
		\textbf{OBJECTIVE FUNCTION:}
		\[
		MSE(b_1) = \frac{1}{N}\sum_{i=1}^N (y_i - 16.88 - b_1 x_i)^2
		\]
		
		\vspace{0.1cm}
		\textbf{GRADIENT:}
		\[
		\frac{\partial MSE}{\partial b_1} = -\frac{2}{N}\sum_{i=1}^N (y_i - \hat{y}_i) x_i
		\]
		
		\vspace{0.1cm}
		\textbf{ITERATION 0:}
		\begin{itemize}
			\item $b_1 = 0.8$
			\item Gradient = -48.664
			\item Update: $0.8 - 0.005(-48.664) = 1.043$
		\end{itemize}
		
		\textcolor{myred}{\textbf{EXAM TRAP:}} Initial $b_1 = 0.8$, final $b_1 = 1.059$!
	\end{frame}
	
	\begin{frame}{T56: Gradient Descent - Convergence Example}
		\fontsize{6.5}{7.5}\selectfont
		\textbf{TOPIC: Convergence in Salary Example}
		
		\vspace{0.1cm}
		\textbf{CONVERGENCE TABLE:}
		\begin{tabular}{|c|c|c|}
			\hline
			Iteration & $b_1$ & Gradient \\
			\hline
			0 & 0.800 & -48.664 \\
			1 & 1.043 & -3.024 \\
			2 & 1.058 & -0.188 \\
			3 & 1.059 & -0.012 \\
			4 & 1.059 & -0.001 \\
			5 & 1.059 & 0.000 \\
			\hline
		\end{tabular}
		
		\vspace{0.1cm}
		\textbf{CONVERGENCE:}
		\begin{itemize}
			\item Gradient approaches zero
			\item $b_1$ stabilizes at 1.059
			\item Only 5 iterations needed
		\end{itemize}
		
		\vspace{0.1cm}
		\textbf{COMPARISON:}
		\begin{itemize}
			\item OLS solution: $\hat{b}_1 = 1.06$
			\item Gradient descent: 1.059 (approximation)
		\end{itemize}
		
		\textcolor{myred}{\textbf{EXAM TRAP:}} Converged in \textcolor{myblue}{5 iterations} in the example!
	\end{frame}
	
	\begin{frame}{T57: Gradient Descent - Epoch Definition}
		\fontsize{6.5}{7.5}\selectfont
		\textbf{TOPIC: Epoch in Machine Learning}
		
		\vspace{0.1cm}
		\textbf{DEFINITION:}
		\begin{itemize}
			\item One iteration including calculating loss and adjusting weights using whole training sample
		\end{itemize}
		
		\vspace{0.1cm}
		\textbf{FOR BATCH GRADIENT DESCENT:}
		\begin{itemize}
			\item 1 epoch = 1 weight update
		\end{itemize}
		
		\vspace{0.1cm}
		\textbf{FOR STOCHASTIC GRADIENT DESCENT:}
		\begin{itemize}
			\item 1 epoch = N weight updates
		\end{itemize}
		
		\vspace{0.1cm}
		\textbf{EXAMPLE:}
		\begin{itemize}
			\item Training data: 1000 observations
			\item Batch GD: 1 update per epoch
			\item SGD: 1000 updates per epoch
		\end{itemize}
		
		\textcolor{myred}{\textbf{EXAM TRAP:}} Epoch = \textcolor{myblue}{one full pass} through training data!
	\end{frame}
	
	\begin{frame}{T58: Gradient Descent - Local Minima vs Global}
		\fontsize{6.5}{7.5}\selectfont
		\textbf{TOPIC: Local vs Global Minimum}
		
		\vspace{0.1cm}
		\textbf{LOCAL MINIMUM:}
		\begin{itemize}
			\item Gradient = 0
			\item Lower than immediate neighbors
			\item May not be lowest overall
		\end{itemize}
		
		\vspace{0.1cm}
		\textbf{GLOBAL MINIMUM:}
		\begin{itemize}
			\item Lowest point of entire function
			\item What we ideally want
		\end{itemize}
		
		\vspace{0.1cm}
		\textbf{CHALLENGE:}
		\begin{itemize}
			\item Gradient descent can get stuck in local minima
			\item No guarantee of finding global minimum
		\end{itemize}
		
		\vspace{0.1cm}
		\textbf{SOLUTIONS:}
		\begin{itemize}
			\item Momentum
			\item Multiple random starts
			\item Stochastic gradient descent
		\end{itemize}
		
		\textcolor{myred}{\textbf{EXAM TRAP:}} Gradient descent may find \textcolor{myblue}{local, not global}, minimum!
	\end{frame}
	
	\begin{frame}{T59: Gradient Descent - Plateau}
		\fontsize{6.5}{7.5}\selectfont
		\textbf{TOPIC: Plateaus in Gradient Descent}
		
		\vspace{0.1cm}
		\textbf{DEFINITION:}
		\begin{itemize}
			\item Flat region where gradient is near zero
			\item Not a minimum
		\end{itemize}
		
		\vspace{0.1cm}
		\textbf{PROBLEM:}
		\begin{itemize}
			\item Gradient descent slows down
			\item May appear converged but isn't
		\end{itemize}
		
		\vspace{0.1cm}
		\textbf{SOLUTION:}
		\begin{itemize}
			\item Momentum helps
			\item Adaptive learning rates
		\end{itemize}
		
		\vspace{0.1cm}
		\textbf{EXAMPLE:}
		\begin{itemize}
			\item Flat region in loss landscape
			\item Gradient small but not zero
			\item Updates tiny
		\end{itemize}
		
		\textcolor{myred}{\textbf{EXAM TRAP:}} Plateaus \textcolor{myblue}{slow down} gradient descent!
	\end{frame}
	
	\begin{frame}{T60: Gradient Descent - Summary}
		\fontsize{6.5}{7.5}\selectfont
		\textbf{TOPIC: Key Advantages of Gradient Descent}
		
		\vspace{0.1cm}
		\textbf{KEY ADVANTAGE:}
		\begin{itemize}
			\item Can optimize complex, high-dimensional, nonlinear models
			\item Handles complex models
			\item Works with high dimensions
			\item No closed-form needed
		\end{itemize}
		
		\vspace{0.1cm}
		\textbf{APPLICATIONS:}
		\begin{itemize}
			\item Neural networks
			\item Logistic regression
			\item Any differentiable model
		\end{itemize}
		
		\vspace{0.1cm}
		\textbf{LIMITATIONS:}
		\begin{itemize}
			\item Can get stuck in local minima
			\item Sensitive to learning rate
			\item May converge slowly
		\end{itemize}
		
		\textcolor{myred}{\textbf{EXAM TRAP:}} Gradient descent = \textcolor{myblue}{workhorse of modern ML}!
	\end{frame}
	
	% ============================================================
	% SECTION 5: BACKPROPAGATION
	% ============================================================
	
	\begin{frame}{T61: Backpropagation Definition}
		\fontsize{6.5}{7.5}\selectfont
		\textbf{TOPIC: Backpropagation - Definition and Overview}
		
		\vspace{0.1cm}
		\textbf{DEFINITION:}
		\begin{itemize}
			\item Algorithm that computes gradients of loss function with respect to all weights
			\item Uses chain rule of differentiation
			\item Foundation of neural network training
		\end{itemize}
		
		\vspace{0.1cm}
		\textbf{PROCESS:}
		\begin{enumerate}
			\item Forward pass: compute output and loss
			\item Backward pass: compute gradients
			\item Update weights using gradient descent
		\end{enumerate}
		
		\vspace{0.1cm}
		\textbf{KEY INSIGHT:}
		\begin{itemize}
			\item Consider each layer separately
			\item Apply chain rule at each step
			\item Propagate errors backward
		\end{itemize}
		
		\vspace{0.1cm}
		\textbf{WHY "BACK":}
		\begin{itemize}
			\item Start at output layer
			\item Work backward through layers
		\end{itemize}
		
		\textcolor{myred}{\textbf{EXAM TRAP:}} Backpropagation uses the \textcolor{myblue}{chain rule}!
	\end{frame}
	
	\begin{frame}{T62: Backpropagation - Chain Rule}
		\fontsize{6.5}{7.5}\selectfont
		\textbf{TOPIC: Chain Rule Foundation of Backpropagation}
		
		\vspace{0.1cm}
		\textbf{CHAIN RULE:}
		\[
		\frac{dy}{dx} = \frac{dy}{du} \cdot \frac{du}{dx}
		\]
		
		\vspace{0.1cm}
		\textbf{NEURAL NETWORK APPLICATION:}
		\begin{itemize}
			\item Output is function of function
			\item Multiple layers compose functions
			\item Chain rule computes gradients through composition
		\end{itemize}
		
		\vspace{0.1cm}
		\textbf{EXAMPLE:}
		\begin{itemize}
			\item $L = f(g(h(x)))$
			\item $\frac{\partial L}{\partial x} = f'(g) \cdot g'(h) \cdot h'(x)$
		\end{itemize}
		
		\vspace{0.1cm}
		\textbf{WHY IT MATTERS:}
		\begin{itemize}
			\item Enables gradient computation for deep networks
			\item Breaks complex problem into simpler steps
			\item Foundation of automatic differentiation
		\end{itemize}
		
		\textcolor{myred}{\textbf{EXAM TRAP:}} Chain rule is the \textcolor{myblue}{foundation} of backpropagation!
	\end{frame}
	
	\begin{frame}{T63: Backpropagation - Forward Pass}
		\fontsize{6.5}{7.5}\selectfont
		\textbf{TOPIC: Forward Pass in Backpropagation}
		
		\vspace{0.1cm}
		\textbf{DEFINITION:}
		\begin{itemize}
			\item Input data passed through network to compute output and loss
		\end{itemize}
		
		\vspace{0.1cm}
		\textbf{PROCESS:}
		\begin{enumerate}
			\item Input features fed to network
			\item Hidden layer activations computed
			\item Output computed
			\item Loss calculated
		\end{enumerate}
		
		\vspace{0.1cm}
		\textbf{EXAMPLE (HIDDEN LAYER):}
		\[
		H_j = f\left(\sum_i w_{i,j} x_i\right)
		\]
		
		\vspace{0.1cm}
		\textbf{EXAMPLE (OUTPUT LAYER):}
		\[
		H_k = f\left(\sum_j w_{j,k} H_j\right)
		\]
		
		\textcolor{myred}{\textbf{EXAM TRAP:}} Forward pass computes \textcolor{myblue}{output and loss}!
	\end{frame}
	
	\begin{frame}{T64: Backpropagation - Backward Pass}
		\fontsize{6.5}{7.5}\selectfont
		\textbf{TOPIC: Backward Pass in Backpropagation}
		
		\vspace{0.1cm}
		\textbf{DEFINITION:}
		\begin{itemize}
			\item Gradients of loss with respect to weights computed using chain rule
		\end{itemize}
		
		\vspace{0.1cm}
		\textbf{PROCESS:}
		\begin{enumerate}
			\item Start at output layer
			\item Compute error at output
			\item Propagate error backward
			\item Compute gradients for each weight
		\end{enumerate}
		
		\vspace{0.1cm}
		\textbf{EXAMPLE (OUTPUT LAYER):}
		\[
		e_k = -(y - H_k) \cdot f'(z_k)
		\]
		
		\vspace{0.1cm}
		\textbf{EXAMPLE (HIDDEN LAYER):}
		\[
		e_j = e_k \cdot w_{j,k} \cdot f'(z_j)
		\]
		
		\textcolor{myred}{\textbf{EXAM TRAP:}} Backward pass computes \textcolor{myblue}{gradients}!
	\end{frame}
	
	\begin{frame}{T65: Backpropagation - Weight Update}
		\fontsize{6.5}{7.5}\selectfont
		\textbf{TOPIC: Weight Update in Backpropagation}
		
		\vspace{0.1cm}
		\textbf{UPDATE RULE:}
		\[
		w_{j,k}^{new} = w_{j,k}^{old} - \eta \frac{\partial \mathcal{L}}{\partial w_{j,k}}
		\]
		
		\vspace{0.1cm}
		\textbf{USING ERRORS:}
		\[
		w_{j,k}^{new} = w_{j,k}^{old} - \eta \cdot e_k \cdot H_j
		\]
		
		\vspace{0.1cm}
		\textbf{PROCESS:}
		\begin{enumerate}
			\item Compute gradient
			\item Multiply by learning rate
			\item Subtract from current weight
		\end{enumerate}
		
		\vspace{0.1cm}
		\textbf{EXAMPLE:}
		\begin{itemize}
			\item $w = 0.01$, $\eta = 0.1$, $e_k = -15.44$, $H_j = 0.5$
			\item $w^{new} = 0.01 - 0.1(-15.44)(0.5) = 0.782$
		\end{itemize}
		
		\textcolor{myred}{\textbf{EXAM TRAP:}} Weights updated by \textcolor{myblue}{subtracting} gradient $\times$ learning rate!
	\end{frame}
	
	\begin{frame}{T66: Backpropagation - Activation Function Derivative}
		\fontsize{6.5}{7.5}\selectfont
		\textbf{TOPIC: Sigmoid Activation Function Derivative}
		
		\vspace{0.1cm}
		\textbf{SIGMOID FUNCTION:}
		\[
		f(z) = \frac{1}{1 + e^{-z}}
		\]
		
		\vspace{0.1cm}
		\textbf{DERIVATIVE:}
		\[
		f'(z) = f(z)(1 - f(z))
		\]
		
		\vspace{0.1cm}
		\textbf{DERIVATION:}
		\begin{itemize}
			\item $f(z) = (1 + e^{-z})^{-1}$
			\item $f'(z) = -1(1 + e^{-z})^{-2} \cdot (-e^{-z})$
			\item $= \frac{e^{-z}}{(1 + e^{-z})^2}$
			\item $= f(z)(1 - f(z))$
		\end{itemize}
		
		\vspace{0.1cm}
		\textbf{IMPORTANCE:}
		\begin{itemize}
			\item Used in backpropagation
			\item Easy to compute from $f(z)$
		\end{itemize}
		
		\textcolor{myred}{\textbf{EXAM TRAP:}} Sigmoid derivative = \textcolor{myblue}{$f(z)(1-f(z))$}!
	\end{frame}
	
	\begin{frame}{T67: Backpropagation - Output Layer Error}
		\fontsize{6.5}{7.5}\selectfont
		\textbf{TOPIC: Output Layer Error Calculation}
		
		\vspace{0.1cm}
		\textbf{FORMULA:}
		\[
		e_k = -(y - H_k) \cdot f'(z_k)
		\]
		
		\vspace{0.1cm}
		\textbf{COMPONENTS:}
		\begin{itemize}
			\item $y$: Actual target
			\item $H_k$: Predicted output
			\item $f'(z_k)$: Activation derivative
		\end{itemize}
		
		\vspace{0.1cm}
		\textbf{FOR LINEAR ACTIVATION:}
		\begin{itemize}
			\item $f'(z_k) = 1$
			\item $e_k = -(y - H_k) = H_k - y$
		\end{itemize}
		
		\vspace{0.1cm}
		\textbf{EXAMPLE:}
		\begin{itemize}
			\item $y = 15.46$, $H_k = 0.02$
			\item $e_k = -(15.46 - 0.02) = -15.44$
		\end{itemize}
		
		\textcolor{myred}{\textbf{EXAM TRAP:}} Output error includes \textcolor{myblue}{activation derivative}!
	\end{frame}
	
	\begin{frame}{T68: Backpropagation - Hidden Layer Error}
		\fontsize{6.5}{7.5}\selectfont
		\textbf{TOPIC: Hidden Layer Error Calculation}
		
		\vspace{0.1cm}
		\textbf{FORMULA:}
		\[
		e_j = e_k \cdot w_{j,k} \cdot f'(z_j)
		\]
		
		\vspace{0.1cm}
		\textbf{COMPONENTS:}
		\begin{itemize}
			\item $e_k$: Error from next layer
			\item $w_{j,k}$: Weight connecting to next layer
			\item $f'(z_j)$: Activation derivative
		\end{itemize}
		
		\vspace{0.1cm}
		\textbf{EXAMPLE:}
		\begin{itemize}
			\item $e_k = -15.44$, $w = 0.01$
			\item $f'(z_j) = 0.5 \times 0.5 = 0.25$
			\item $e_j = -15.44 \times 0.01 \times 0.25 = -0.0386$
		\end{itemize}
		
		\vspace{0.1cm}
		\textbf{PROCESS:}
		\begin{itemize}
			\item Start with output error
			\item Work backward through layers
			\item Multiply by weights and derivatives
		\end{itemize}
		
		\textcolor{myred}{\textbf{EXAM TRAP:}} Hidden error uses \textcolor{myblue}{weights and activation derivative}!
	\end{frame}
	
	\begin{frame}{T69: Backpropagation - Summary of Steps}
		\fontsize{6.5}{7.5}\selectfont
		\textbf{TOPIC: Summary of Backpropagation Steps}
		
		\vspace{0.1cm}
		\textbf{CORRECT ORDER:}
		\begin{enumerate}
			\item Forward pass: compute output
			\item Error calculation: compare to actual
			\item Backward pass: compute gradients
			\item Update weights using gradient descent
		\end{enumerate}
		
		\vspace{0.1cm}
		\textbf{REPEAT:}
		\begin{itemize}
			\item For each data point
			\item Until convergence
		\end{itemize}
		
		\vspace{0.1cm}
		\textbf{CONVERGENCE CRITERIA:}
		\begin{itemize}
			\item Loss below threshold
			\item Improvement below threshold
			\item Maximum iterations reached
			\item Validation error increases
		\end{itemize}
		
		\textcolor{myred}{\textbf{EXAM TRAP:}} Forward $\to$ Error $\to$ Backward $\to$ Update!
	\end{frame}
	
	\begin{frame}{T70: Backpropagation - Convergence Check}
		\fontsize{6.5}{7.5}\selectfont
		\textbf{TOPIC: When Does Backpropagation Stop?}
		
		\vspace{0.1cm}
		\textbf{CONVERGENCE CRITERIA:}
		\begin{enumerate}
			\item Loss below threshold
			\item Improvement below threshold
			\item Maximum iterations reached
			\item Validation error increases
		\end{enumerate}
		
		\vspace{0.1cm}
		\textbf{EPOCHS:}
		\begin{itemize}
			\item Multiple passes through data
			\item Number depends on problem
		\end{itemize}
		
		\vspace{0.1cm}
		\textbf{EXAMPLE:}
		\begin{itemize}
			\item Stop when loss < 0.001
			\item Stop after 100 epochs
			\item Stop when validation error increases
		\end{itemize}
		
		\textcolor{myred}{\textbf{EXAM TRAP:}} Multiple stopping criteria, not just one!
	\end{frame}
	
	\begin{frame}{T71: Backpropagation - Vanishing Gradient}
		\fontsize{6.5}{7.5}\selectfont
		\textbf{TOPIC: Vanishing Gradient Problem}
		
		\vspace{0.1cm}
		\textbf{DEFINITION:}
		\begin{itemize}
			\item Gradients become very small during backpropagation
			\item Especially in early layers
			\item Product of many small derivatives
		\end{itemize}
		
		\vspace{0.1cm}
		\textbf{CAUSE:}
		\begin{itemize}
			\item Chain rule multiplies derivatives
			\item Sigmoid derivatives are $<1$
			\item Product quickly approaches zero
		\end{itemize}
		
		\vspace{0.1cm}
		\textbf{SOLUTION:}
		\begin{itemize}
			\item Use ReLU activation
			\item Batch normalization
			\item LSTM architectures
		\end{itemize}
		
		\vspace{0.1cm}
		\textbf{CONSEQUENCE:}
		\begin{itemize}
			\item Early layers don't learn
			\item Training very slow
		\end{itemize}
		
		\textcolor{myred}{\textbf{EXAM TRAP:}} Vanishing gradient = \textcolor{myblue}{very small gradients} in early layers!
	\end{frame}
	
	\begin{frame}{T72: Backpropagation - Exploding Gradient}
		\fontsize{6.5}{7.5}\selectfont
		\textbf{TOPIC: Exploding Gradient Problem}
		
		\vspace{0.1cm}
		\textbf{DEFINITION:}
		\begin{itemize}
			\item Gradients become very large during backpropagation
			\item Product of many large derivatives
			\item Unstable parameter updates
		\end{itemize}
		
		\vspace{0.1cm}
		\textbf{CONSEQUENCES:}
		\begin{itemize}
			\item Loss diverges
			\item Training fails
		\end{itemize}
		
		\vspace{0.1cm}
		\textbf{SOLUTION:}
		\begin{itemize}
			\item Gradient clipping
			\item Lower learning rate
			\item Batch normalization
		\end{itemize}
		
		\vspace{0.1cm}
		\textbf{OPPOSITE PROBLEM:}
		\begin{itemize}
			\item Vanishing gradient (very small)
		\end{itemize}
		
		\textcolor{myred}{\textbf{EXAM TRAP:}} Exploding gradient = \textcolor{myblue}{very large gradients}!
	\end{frame}
	
	\begin{frame}{T73: Backpropagation - ReLU Activation}
		\fontsize{6.5}{7.5}\selectfont
		\textbf{TOPIC: ReLU Activation Function}
		
		\vspace{0.1cm}
		\textbf{RELU FUNCTION:}
		\[
		f(z) = \max(0, z)
		\]
		
		\vspace{0.1cm}
		\textbf{DERIVATIVE:}
		\[
		f'(z) = \begin{cases} 1 & \text{if } z > 0 \\ 0 & \text{if } z \leq 0 \end{cases}
		\]
		
		\vspace{0.1cm}
		\textbf{ADVANTAGES:}
		\begin{itemize}
			\item Derivative is 1 for positive values
			\item No vanishing gradient for positive values
			\item Computationally efficient
		\end{itemize}
		
		\vspace{0.1cm}
		\textbf{WHY IT HELPS:}
		\begin{itemize}
			\item Avoids vanishing gradient problem
			\item Faster training
			\item Standard in modern networks
		\end{itemize}
		
		\textcolor{myred}{\textbf{EXAM TRAP:}} ReLU \textcolor{myblue}{avoids vanishing gradient}!
	\end{frame}
	
	\begin{frame}{T74: Backpropagation - Batch Normalization}
		\fontsize{6.5}{7.5}\selectfont
		\textbf{TOPIC: Batch Normalization}
		
		\vspace{0.1cm}
		\textbf{DEFINITION:}
		\begin{itemize}
			\item Adds normalization layers between hidden layers
			\item Normalizes layer inputs
		\end{itemize}
		
		\vspace{0.1cm}
		\textbf{PROCESS:}
		\begin{itemize}
			\item Subtract mean
			\item Divide by standard deviation
			\item Apply between hidden layers
		\end{itemize}
		
		\vspace{0.1cm}
		\textbf{BENEFITS:}
		\begin{itemize}
			\item Reduces internal covariate shift
			\item Helps with vanishing/exploding gradients
			\item Allows higher learning rates
		\end{itemize}
		
		\vspace{0.1cm}
		\textbf{ANALOGY:}
		\begin{itemize}
			\item Like normalizing input data
			\item But applied to hidden layers
		\end{itemize}
		
		\textcolor{myred}{\textbf{EXAM TRAP:}} Batch normalization helps with \textcolor{myblue}{gradient problems}!
	\end{frame}
	
	\begin{frame}{T75: Backpropagation - Summary}
		\fontsize{6.5}{7.5}\selectfont
		\textbf{TOPIC: Summary of Backpropagation}
		
		\vspace{0.1cm}
		\textbf{TRUE STATEMENTS:}
		\begin{itemize}
			\item Computes gradients for all weights
			\item Uses chain rule
			\item Requires storing activations
		\end{itemize}
		
		\vspace{0.1cm}
		\textbf{FALSE STATEMENT:}
		\begin{itemize}
			\item Backpropagation does NOT guarantee global minimum
			\item Can get stuck in local minima
			\item Can suffer from vanishing/exploding gradients
		\end{itemize}
		
		\vspace{0.1cm}
		\textbf{KEY POINTS:}
		\begin{itemize}
			\item Foundation of neural network training
			\item Chain rule is essential
			\item Multiple computational challenges exist
		\end{itemize}
		
		\textcolor{myred}{\textbf{EXAM TRAP:}} Backpropagation does \textcolor{myred}{NOT} guarantee global minimum!
	\end{frame}
	
	% ============================================================
	% SECTION 6: BIAS-VARIANCE TRADEOFF
	% ============================================================
	
	\begin{frame}{T76: Bias-Variance Tradeoff Definition}
		\fontsize{6.5}{7.5}\selectfont
		\textbf{TOPIC: Bias-Variance Tradeoff - Definition}
		
		\vspace{0.1cm}
		\textbf{DEFINITION:}
		\begin{itemize}
			\item Tradeoff between choosing a model that is not complex enough (underfitting) and one that is too complex (overfitting)
			\item Balances bias and variance
		\end{itemize}
		
		\vspace{0.1cm}
		\textbf{TOTAL ERROR:}
		\[
		\text{Total Error} = \text{Bias}^2 + \text{Variance} + \text{Irreducible Error}
		\]
		
		\vspace{0.1cm}
		\textbf{GOAL:}
		\begin{itemize}
			\item Find optimal complexity
			\item Balance bias and variance
		\end{itemize}
		
		\vspace{0.1cm}
		\textbf{VISUAL:}
		\begin{itemize}
			\item Simple model: high bias, low variance
			\item Complex model: low bias, high variance
			\item Optimal: balance
		\end{itemize}
		
		\textcolor{myred}{\textbf{EXAM TRAP:}} Bias = \textcolor{myblue}{underfitting}; Variance = \textcolor{myblue}{overfitting}!
	\end{frame}
	
	\begin{frame}{T77: Bias Definition}
		\fontsize{6.5}{7.5}\selectfont
		\textbf{TOPIC: Bias in the Bias-Variance Tradeoff}
		
		\vspace{0.1cm}
		\textbf{DEFINITION:}
		\begin{itemize}
			\item Error from model's assumptions and limitations
			\item High bias = underfitting
			\item Model too simple
		\end{itemize}
		
		\vspace{0.1cm}
		\textbf{EXAMPLE:}
		\begin{itemize}
			\item Linear model for nonlinear data
			\item Cannot capture true relationship
		\end{itemize}
		
		\vspace{0.1cm}
		\textbf{CHARACTERISTICS:}
		\begin{itemize}
			\item Poor training performance
			\item Poor test performance
		\end{itemize}
		
		\vspace{0.1cm}
		\textbf{WHY IT MATTERS:}
		\begin{itemize}
			\item Model systematically wrong
			\item Predictions biased
		\end{itemize}
		
		\textcolor{myred}{\textbf{EXAM TRAP:}} Bias = \textcolor{myblue}{error from assumptions} (underfitting)!
	\end{frame}
	
	\begin{frame}{T78: Variance Definition}
		\fontsize{6.5}{7.5}\selectfont
		\textbf{TOPIC: Variance in the Bias-Variance Tradeoff}
		
		\vspace{0.1cm}
		\textbf{DEFINITION:}
		\begin{itemize}
			\item Error from sensitivity to fluctuations in training data
			\item High variance = overfitting
			\item Model too complex
		\end{itemize}
		
		\vspace{0.1cm}
		\textbf{EXAMPLE:}
		\begin{itemize}
			\item High-degree polynomial
			\item Fits noise in training data
		\end{itemize}
		
		\vspace{0.1cm}
		\textbf{CHARACTERISTICS:}
		\begin{itemize}
			\item Good training performance
			\item Poor test performance
		\end{itemize}
		
		\vspace{0.1cm}
		\textbf{WHY IT MATTERS:}
		\begin{itemize}
			\item Model unstable
			\item Small data changes → big prediction changes
		\end{itemize}
		
		\textcolor{myred}{\textbf{EXAM TRAP:}} Variance = \textcolor{myblue}{error from sensitivity} (overfitting)!
	\end{frame}
	
	\begin{frame}{T79: Underfitting Definition}
		\fontsize{6.5}{7.5}\selectfont
		\textbf{TOPIC: Underfitting - Definition and Causes}
		
		\vspace{0.1cm}
		\textbf{DEFINITION:}
		\begin{itemize}
			\item Model too simple to capture underlying patterns
			\item High bias
			\item Fails to capture relevant patterns
		\end{itemize}
		
		\vspace{0.1cm}
		\textbf{EXAMPLE:}
		\begin{itemize}
			\item Linear model for quadratic relationship
			\item Missing interaction terms
		\end{itemize}
		
		\vspace{0.1cm}
		\textbf{SYMPTOMS:}
		\begin{itemize}
			\item Poor training performance
			\item Poor test performance
		\end{itemize}
		
		\vspace{0.1cm}
		\textbf{CAUSES:}
		\begin{itemize}
			\item Too few features
			\item Model too simple
			\item Excessive regularization
		\end{itemize}
		
		\textcolor{myred}{\textbf{EXAM TRAP:}} Underfitting = \textcolor{myblue}{model too simple}!
	\end{frame}
	
	\begin{frame}{T80: Overfitting Definition}
		\fontsize{6.5}{7.5}\selectfont
		\textbf{TOPIC: Overfitting - Definition and Causes}
		
		\vspace{0.1cm}
		\textbf{DEFINITION:}
		\begin{itemize}
			\item Model too complex, fits noise in training data
			\item High variance
			\item Learns random noise
		\end{itemize}
		
		\vspace{0.1cm}
		\textbf{EXAMPLE:}
		\begin{itemize}
			\item High-degree polynomial
			\item Neural network with too many parameters
		\end{itemize}
		
		\vspace{0.1cm}
		\textbf{SYMPTOMS:}
		\begin{itemize}
			\item Good training performance
			\item Poor test performance
		\end{itemize}
		
		\vspace{0.1cm}
		\textbf{CAUSES:}
		\begin{itemize}
			\item Too many parameters
			\item Too little data
			\item No regularization
		\end{itemize}
		
		\textcolor{myred}{\textbf{EXAM TRAP:}} Overfitting = \textcolor{myblue}{model too complex}!
	\end{frame}
	
	\begin{frame}{T81: Bias-Variance - Model Complexity}
		\fontsize{6.5}{7.5}\selectfont
		\textbf{TOPIC: Effect of Model Complexity on Bias and Variance}
		
		\vspace{0.1cm}
		\textbf{AS COMPLEXITY INCREASES:}
		\begin{itemize}
			\item Bias decreases (fits training data better)
			\item Variance increases (more sensitive to training data)
		\end{itemize}
		
		\vspace{0.1cm}
		\textbf{VISUAL:}
		\begin{itemize}
			\item Simple model: high bias, low variance
			\item Complex model: low bias, high variance
			\item Optimal: balance
		\end{itemize}
		
		\vspace{0.1cm}
		\textbf{EXAMPLE:}
		\begin{itemize}
			\item Linear regression: high bias, low variance
			\item 20th degree polynomial: low bias, high variance
		\end{itemize}
		
		\vspace{0.1cm}
		\textbf{KEY INSIGHT:}
		\begin{itemize}
			\item Tradeoff is fundamental
			\item Cannot minimize both simultaneously
		\end{itemize}
		
		\textcolor{myred}{\textbf{EXAM TRAP:}} More complexity = \textcolor{myblue}{lower bias, higher variance}!
	\end{frame}
	
	\begin{frame}{T82: Optimal Model Complexity}
		\fontsize{6.5}{7.5}\selectfont
		\textbf{TOPIC: Finding Optimal Model Complexity}
		
		\vspace{0.1cm}
		\textbf{DEFINITION:}
		\begin{itemize}
			\item Complexity that minimizes total error
			\item Total error = Bias² + Variance + Irreducible Error
		\end{itemize}
		
		\vspace{0.1cm}
		\textbf{TRADEOFF:}
		\begin{itemize}
			\item Too simple: high bias
			\item Too complex: high variance
			\item Optimal: balance
		\end{itemize}
		
		\vspace{0.1cm}
		\textbf{HOW TO FIND:}
		\begin{itemize}
			\item Cross-validation
			\item Validation set performance
			\item Regularization
		\end{itemize}
		
		\vspace{0.1cm}
		\textbf{EXAMPLE:}
		\begin{itemize}
			\item Quadratic model often optimal
			\item Balances underfitting and overfitting
		\end{itemize}
		
		\textcolor{myred}{\textbf{EXAM TRAP:}} Optimal = \textcolor{myblue}{minimum total error}!
	\end{frame}
	
	\begin{frame}{T83: House Price Example - Underfitting}
		\fontsize{6.5}{7.5}\selectfont
		\textbf{TOPIC: Underfitting in House Price Example}
		
		\vspace{0.1cm}
		\textbf{LINEAR REGRESSION:}
		\begin{itemize}
			\item Forces strictly decreasing relationship
			\item Cannot capture non-monotonic pattern
			\item Underfits the data
		\end{itemize}
		
		\vspace{0.1cm}
		\textbf{EVIDENCE:}
		\begin{itemize}
			\item Training MSE: 166.67
			\item Test MSE: 161.12
			\item Both high
		\end{itemize}
		
		\vspace{0.1cm}
		\textbf{WHY:}
		\begin{itemize}
			\item House prices non-monotonic in age
			\item Both very new and very old houses expensive
			\item Linear model can't capture U-shape
		\end{itemize}
		
		\textcolor{myred}{\textbf{EXAM TRAP:}} Linear regression \textcolor{myblue}{underfits} house price data!
	\end{frame}
	
	\begin{frame}{T84: House Price Example - Overfitting}
		\fontsize{6.5}{7.5}\selectfont
		\textbf{TOPIC: Overfitting in House Price Example}
		
		\vspace{0.1cm}
		\textbf{NINTH-DEGREE POLYNOMIAL:}
		\begin{itemize}
			\item Captures data-specific noise
			\item Curved shape between 20-30 years
			\item Poor generalization
		\end{itemize}
		
		\vspace{0.1cm}
		\textbf{EVIDENCE:}
		\begin{itemize}
			\item Training MSE: 122.40 (lowest)
			\item Test MSE: 139.39 (higher than quadratic)
		\end{itemize}
		
		\vspace{0.1cm}
		\textbf{LESSON:}
		\begin{itemize}
			\item Complex models can overfit
			\item Test performance reveals overfitting
		\end{itemize}
		
		\textcolor{myred}{\textbf{EXAM TRAP:}} Ninth-degree polynomial \textcolor{myblue}{overfits}!
	\end{frame}
	
	\begin{frame}{T85: House Price Example - Optimal}
		\fontsize{6.5}{7.5}\selectfont
		\textbf{TOPIC: Optimal Model in House Price Example}
		
		\vspace{0.1cm}
		\textbf{QUADRATIC REGRESSION:}
		\begin{itemize}
			\item Captures U-shape pattern
			\item Training MSE: 131.81
			\item Test MSE: 136.57 (lowest)
		\end{itemize}
		
		\vspace{0.1cm}
		\textbf{WHY BEST:}
		\begin{itemize}
			\item Balanced complexity
			\item Generalizes well
			\item Captures true relationship
		\end{itemize}
		
		\vspace{0.1cm}
		\textbf{COMPARISON:}
		\begin{itemize}
			\item Linear: underfits (high bias)
			\item Quadratic: optimal balance
			\item Ninth-degree: overfits (high variance)
		\end{itemize}
		
		\textcolor{myred}{\textbf{EXAM TRAP:}} Quadratic strikes \textcolor{myblue}{best balance}!
	\end{frame}
	
	% ============================================================
	% SECTION 7: REGULARIZATION
	% ============================================================
	
	\begin{frame}{T86: Regularization Definition}
		\fontsize{6.5}{7.5}\selectfont
		\textbf{TOPIC: Regularization - Definition and Purpose}
		
		\vspace{0.1cm}
		\textbf{DEFINITION:}
		\begin{itemize}
			\item Technique to reduce overfitting by adding penalty term to loss function
			\item Shrinks coefficients toward zero
			\item Prevents overfitting
		\end{itemize}
		
		\vspace{0.1cm}
		\textbf{LOSS FUNCTION:}
		\[
		L = \text{RSS} + \lambda \cdot \text{Penalty}
		\]
		
		\vspace{0.1cm}
		\textbf{TYPES:}
		\begin{itemize}
			\item \textcolor{myblue}{Ridge (L2):} $\lambda \sum \beta_j^2$
			\item \textcolor{myblue}{LASSO (L1):} $\lambda \sum |\beta_j|$
			\item \textcolor{myblue}{Elastic Net:} both
		\end{itemize}
		
		\vspace{0.1cm}
		\textbf{BENEFITS:}
		\begin{itemize}
			\item Simplifies models
			\item Reduces overfitting
			\item Improves generalization
		\end{itemize}
		
		\textcolor{myred}{\textbf{EXAM TRAP:}} Regularization adds a \textcolor{myblue}{penalty term}!
	\end{frame}
	
	\begin{frame}{T87: Ridge Regression}
		\fontsize{6.5}{7.5}\selectfont
		\textbf{TOPIC: Ridge Regression (L2 Regularization)}
		
		\vspace{0.1cm}
		\textbf{LOSS FUNCTION:}
		\[
		L = \frac{1}{N}\sum(y_i - \hat{y}_i)^2 + \lambda \sum_{j=1}^m \hat{\beta}_j^2
		\]
		
		\vspace{0.1cm}
		\textbf{CHARACTERISTICS:}
		\begin{itemize}
			\item Shrinks coefficients
			\item Does NOT set to zero
			\item Handles multicollinearity
		\end{itemize}
		
		\vspace{0.1cm}
		\textbf{WHY L2:}
		\begin{itemize}
			\item Squared penalty term
			\item Differentiable everywhere
			\item Analytic solution exists
		\end{itemize}
		
		\vspace{0.1cm}
		\textbf{EXAMPLE:}
		\begin{itemize}
			\item OLS coefficient: -23.22
			\item Ridge ($\lambda=0.1$): -6.55
			\item Ridge ($\lambda=0.5$): -2.00
		\end{itemize}
		
		\textcolor{myred}{\textbf{EXAM TRAP:}} Ridge \textcolor{myblue}{shrinks but doesn't zero} coefficients!
	\end{frame}
	
	\begin{frame}{T88: LASSO Regression}
		\fontsize{6.5}{7.5}\selectfont
		\textbf{TOPIC: LASSO Regression (L1 Regularization)}
		
		\vspace{0.1cm}
		\textbf{LOSS FUNCTION:}
		\[
		L = \frac{1}{N}\sum(y_i - \hat{y}_i)^2 + \lambda \sum_{j=1}^m |\hat{\beta}_j|
		\]
		
		\vspace{0.1cm}
		\textbf{CHARACTERISTICS:}
		\begin{itemize}
			\item Sets some coefficients to zero
			\item Performs feature selection
			\item Numerical optimization needed
		\end{itemize}
		
		\vspace{0.1cm}
		\textbf{WHY L1:}
		\begin{itemize}
			\item Absolute value penalty
			\item Creates sparsity
			\item Not differentiable at zero
		\end{itemize}
		
		\vspace{0.1cm}
		\textbf{EXAMPLE:}
		\begin{itemize}
			\item $\lambda = 0.01$: Some coefficients zero
			\item $\lambda = 0.1$: Only 1 non-zero
		\end{itemize}
		
		\textcolor{myred}{\textbf{EXAM TRAP:}} LASSO \textcolor{myblue}{sets coefficients to zero}!
	\end{frame}
	
	\begin{frame}{T89: Elastic Net}
		\fontsize{6.5}{7.5}\selectfont
		\textbf{TOPIC: Elastic Net Regularization}
		
		\vspace{0.1cm}
		\textbf{DEFINITION:}
		\begin{itemize}
			\item Combination of L1 and L2
			\item $\lambda_1 \sum \beta_j^2 + \lambda_2 \sum |\beta_j|$
		\end{itemize}
		
		\vspace{0.1cm}
		\textbf{LOSS FUNCTION:}
		\[
		L = \text{RSS} + \lambda_1 \sum \beta_j^2 + \lambda_2 \sum |\beta_j|
		\]
		
		\vspace{0.1cm}
		\textbf{BENEFITS:}
		\begin{itemize}
			\item Combines both penalties
			\item Shrinks and selects
			\item More flexible
		\end{itemize}
		
		\vspace{0.1cm}
		\textbf{MODERN PREFERENCE:}
		\begin{itemize}
			\item Shift toward Elastic Net
			\item Away from pure LASSO/Ridge
		\end{itemize}
		
		\textcolor{myred}{\textbf{EXAM TRAP:}} Elastic Net = \textcolor{myblue}{L1 + L2}!
	\end{frame}
	
	\begin{frame}{T90: Regularization Parameter}
		\fontsize{6.5}{7.5}\selectfont
		\textbf{TOPIC: Effect of Regularization Parameter $\lambda$}
		
		\vspace{0.1cm}
		\textbf{AS $\lambda$ INCREASES:}
		\begin{itemize}
			\item Penalty becomes stronger
			\item Coefficients shrink more
			\item Model becomes simpler
		\end{itemize}
		
		\vspace{0.1cm}
		\textbf{TRADEOFF:}
		\begin{itemize}
			\item Small $\lambda$: little regularization
			\item Large $\lambda$: heavy regularization
			\item Optimal $\lambda$: balance
		\end{itemize}
		
		\vspace{0.1cm}
		\textbf{EXAMPLE:}
		\begin{itemize}
			\item $\lambda = 0$: OLS
			\item $\lambda = 0.1$: Moderate shrinkage
			\item $\lambda = 10$: Heavy shrinkage
		\end{itemize}
		
		\textcolor{myred}{\textbf{EXAM TRAP:}} Larger $\lambda$ = \textcolor{myblue}{simpler model}!
	\end{frame}
	
	\begin{frame}{T91: LASSO - Feature Selection}
		\fontsize{6.5}{7.5}\selectfont
		\textbf{TOPIC: LASSO as Feature Selection}
		
		\vspace{0.1cm}
		\textbf{LASSO PROPERTY:}
		\begin{itemize}
			\item L1 penalty creates sparsity
			\item Some $\beta_j = 0$ exactly
			\item Those features removed
		\end{itemize}
		
		\vspace{0.1cm}
		\textbf{AS $\lambda$ INCREASES:}
		\begin{itemize}
			\item More coefficients become zero
			\item More features removed
		\end{itemize}
		
		\vspace{0.1cm}
		\textbf{EXAMPLE FROM TEXT:}
		\begin{itemize}
			\item $\lambda = 0.1$: Only 1 non-zero coefficient
		\end{itemize}
		
		\vspace{0.1cm}
		\textbf{WHY USEFUL:}
		\begin{itemize}
			\item Automatic feature selection
			\item Simpler models
			\item Easier interpretation
		\end{itemize}
		
		\textcolor{myred}{\textbf{EXAM TRAP:}} LASSO \textcolor{myblue}{removes features} by zeroing coefficients!
	\end{frame}
	
	\begin{frame}{T92: Ridge vs LASSO}
		\fontsize{6.5}{7.5}\selectfont
		\textbf{TOPIC: Key Differences Between Ridge and LASSO}
		
		\vspace{0.1cm}
		\textbf{RIDGE (L2):}
		\begin{itemize}
			\item Shrinks coefficients
			\item Does NOT set to zero
			\item Handles multicollinearity
		\end{itemize}
		
		\vspace{0.1cm}
		\textbf{LASSO (L1):}
		\begin{itemize}
			\item Can set coefficients to zero
			\item Feature selection
			\item Numerical optimization
		\end{itemize}
		
		\vspace{0.1cm}
		\textbf{CHOICE:}
		\begin{itemize}
			\item Ridge: reduce extreme estimates
			\item LASSO: remove unimportant features
		\end{itemize}
		
		\vspace{0.1cm}
		\textbf{COMPARISON TABLE:}
		\begin{tabular}{|p{2cm}|p{2cm}|p{2cm}|}
			\hline
			& Ridge & LASSO \\
			\hline
			Penalty & L2 & L1 \\
			\hline
			Zeros & No & Yes \\
			\hline
			Feature & No & Yes \\
			\hline
		\end{tabular}
		
		\textcolor{myred}{\textbf{EXAM TRAP:}} Ridge = \textcolor{myblue}{shrink}; LASSO = \textcolor{myblue}{select}!
	\end{frame}
	
	\begin{frame}{T93: Regularization Example}
		\fontsize{6.5}{7.5}\selectfont
		\textbf{TOPIC: Regularization Example - Stock Returns}
		
		\vspace{0.1cm}
		\textbf{OLS RESULTS:}
		\begin{itemize}
			\item Large coefficients
			\item Some opposite signs
			\item Indicates multicollinearity
		\end{itemize}
		
		\vspace{0.1cm}
		\textbf{EXAMPLE:}
		\begin{itemize}
			\item USTB1M: -23.22
			\item USTB3M: +50.64
			\item Opposite signs for correlated features
		\end{itemize}
		
		\vspace{0.1cm}
		\textbf{SOLUTION:}
		\begin{itemize}
			\item Ridge regression
			\item LASSO
		\end{itemize}
		
		\textcolor{myred}{\textbf{EXAM TRAP:}} OLS gave \textcolor{myblue}{large coefficients} due to multicollinearity!
	\end{frame}
	
	\begin{frame}{T94: Ridge Example}
		\fontsize{6.5}{7.5}\selectfont
		\textbf{TOPIC: Ridge Regression Example}
		
		\vspace{0.1cm}
		\textbf{RIDGE RESULTS:}
		\begin{itemize}
			\item Reduced coefficient magnitudes
			\item Higher $\lambda$ = more shrinkage
			\item Some signs changed
		\end{itemize}
		
		\vspace{0.1cm}
		\textbf{EXAMPLE:}
		\begin{itemize}
			\item OLS USTB1M: -23.22
			\item Ridge $\lambda=0.1$: -6.55
			\item Ridge $\lambda=0.5$: -2.00
		\end{itemize}
		
		\vspace{0.1cm}
		\textbf{WHY THIS HELPS:}
		\begin{itemize}
			\item Reduces instability from multicollinearity
			\item More stable predictions
		\end{itemize}
		
		\textcolor{myred}{\textbf{EXAM TRAP:}} Ridge \textcolor{myblue}{reduces coefficient magnitude}!
	\end{frame}
	
	\begin{frame}{T95: LASSO Example}
		\fontsize{6.5}{7.5}\selectfont
		\textbf{TOPIC: LASSO Regression Example}
		
		\vspace{0.1cm}
		\textbf{LASSO RESULTS:}
		\begin{itemize}
			\item Set some coefficients to zero
			\item $\lambda = 0.1$: Only 1 non-zero
			\item Feature selection
		\end{itemize}
		
		\vspace{0.1cm}
		\textbf{EXAMPLE:}
		\begin{itemize}
			\item USTB1M: 0
			\item USTB3M: 0
			\item USTB5Y: -0.71 (only non-zero)
		\end{itemize}
		
		\vspace{0.1cm}
		\textbf{BENEFIT:}
		\begin{itemize}
			\item Simpler model
			\item Easier interpretation
			\item Automatic feature selection
		\end{itemize}
		
		\textcolor{myred}{\textbf{EXAM TRAP:}} LASSO \textcolor{myblue}{sets coefficients to zero}!
	\end{frame}
	
	% ============================================================
	% SECTION 8: CROSS-VALIDATION AND BOOTSTRAPPING
	% ============================================================
	
	\begin{frame}{T96: Cross-Validation Definition}
		\fontsize{6.5}{7.5}\selectfont
		\textbf{TOPIC: K-Fold Cross-Validation}
		
		\vspace{0.1cm}
		\textbf{DEFINITION:}
		\begin{itemize}
			\item Splitting data into k folds
			\item Training on k-1 folds, validating on 1 fold
			\item Repeating k times
		\end{itemize}
		
		\vspace{0.1cm}
		\textbf{K-FOLD CV:}
		\begin{enumerate}
			\item Split data into k folds
			\item For each fold:
			\begin{itemize}
				\item Train on k-1 folds
				\item Validate on 1 fold
			\end{itemize}
			\item Average results
		\end{enumerate}
		
		\vspace{0.1cm}
		\textbf{COMMON VALUES:}
		\begin{itemize}
			\item k = 5 or 10
			\item k = N (LOOCV)
		\end{itemize}
		
		\vspace{0.1cm}
		\textbf{BENEFITS:}
		\begin{itemize}
			\item Efficient use of data
			\item Reliable performance estimate
		\end{itemize}
		
		\textcolor{myred}{\textbf{EXAM TRAP:}} K-fold = train on \textcolor{myblue}{k-1}, validate on \textcolor{myblue}{1}!
	\end{frame}
	
	\begin{frame}{T97: Stratified Cross-Validation}
		\fontsize{6.5}{7.5}\selectfont
		\textbf{TOPIC: Stratified Cross-Validation}
		
		\vspace{0.1cm}
		\textbf{DEFINITION:}
		\begin{itemize}
			\item Preserves class distribution in each fold
			\item Important for imbalanced data
		\end{itemize}
		
		\vspace{0.1cm}
		\textbf{WHY NEEDED:}
		\begin{itemize}
			\item Random folds may miss rare class
			\item Distorted evaluation
		\end{itemize}
		
		\vspace{0.1cm}
		\textbf{EXAMPLE:}
		\begin{itemize}
			\item 100 observations, 10 defaults
			\item 10-fold CV: 9 non-defaults, 1 default per fold
		\end{itemize}
		
		\vspace{0.1cm}
		\textbf{PROCESS:}
		\begin{itemize}
			\item Sample positive and negative separately
			\item Preserve class ratio
		\end{itemize}
		
		\textcolor{myred}{\textbf{EXAM TRAP:}} Stratified CV preserves \textcolor{myblue}{class distribution}!
	\end{frame}
	
	\begin{frame}{T98: Bootstrapping Definition}
		\fontsize{6.5}{7.5}\selectfont
		\textbf{TOPIC: Bootstrapping}
		
		\vspace{0.1cm}
		\textbf{DEFINITION:}
		\begin{itemize}
			\item Sampling with replacement from original data
			\item Creates multiple bootstrap samples
		\end{itemize}
		
		\vspace{0.1cm}
		\textbf{PROCESS:}
		\begin{enumerate}
			\item Create bootstrap sample
			\item Train model
			\item Validate on out-of-bag data
			\item Repeat many times
		\end{enumerate}
		
		\vspace{0.1cm}
		\textbf{KEY FACTS:}
		\begin{itemize}
			\item Some observations repeated
			\item ~37% not sampled (out-of-bag)
			\item Large number of iterations
		\end{itemize}
		
		\vspace{0.1cm}
		\textbf{BENEFITS:}
		\begin{itemize}
			\item Estimates variability
			\item Confidence intervals
		\end{itemize}
		
		\textcolor{myred}{\textbf{EXAM TRAP:}} Bootstrapping samples \textcolor{myblue}{with replacement}!
	\end{frame}
	
	\begin{frame}{T99: Grid Search}
		\fontsize{6.5}{7.5}\selectfont
		\textbf{TOPIC: Grid Search for Hyperparameter Tuning}
		
		\vspace{0.1cm}
		\textbf{DEFINITION:}
		\begin{itemize}
			\item Exhaustive search over predefined hyperparameter grid
			\item Used with cross-validation
		\end{itemize}
		
		\vspace{0.1cm}
		\textbf{PROCESS:}
		\begin{enumerate}
			\item Define grid of hyperparameter values
			\item Evaluate all combinations
			\item Use cross-validation
			\item Select best
		\end{enumerate}
		
		\vspace{0.1cm}
		\textbf{EXAMPLE:}
		\begin{itemize}
			\item $\lambda$ from 0 to 100, step 1
			\item 5-fold CV
			\item 101 × 5 = 505 model fits
		\end{itemize}
		
		\vspace{0.1cm}
		\textbf{DRAWBACK:}
		\begin{itemize}
			\item Computationally expensive
			\item May overfit the grid
		\end{itemize}
		
		\textcolor{myred}{\textbf{EXAM TRAP:}} Grid search is \textcolor{myblue}{exhaustive}!
	\end{frame}
	
	\begin{frame}{T100: Chapter Summary}
		\fontsize{6.5}{7.5}\selectfont
		\textbf{TOPIC: Key Takeaways from Model Estimation}
		
		\vspace{0.1cm}
		\textbf{ESTIMATION METHODS:}
		\begin{itemize}
			\item \textcolor{myblue}{OLS:} Linear models, closed-form
			\item \textcolor{myblue}{NLS:} Nonlinear models, iterative
			\item \textcolor{myblue}{MLE:} Probabilistic, flexible
		\end{itemize}
		
		\vspace{0.1cm}
		\textbf{OPTIMIZATION:}
		\begin{itemize}
			\item \textcolor{myblue}{Gradient Descent:} Negative gradient, learning rate
			\item \textcolor{myblue}{Backpropagation:} Chain rule, neural networks
			\item \textcolor{myblue}{Momentum:} Accelerates convergence
		\end{itemize}
		
		\vspace{0.1cm}
		\textbf{MODEL SELECTION:}
		\begin{itemize}
			\item \textcolor{myblue}{Bias-Variance:} Underfitting vs overfitting
			\item \textcolor{myblue}{Regularization:} Ridge (L2), LASSO (L1), Elastic Net
			\item \textcolor{myblue}{Cross-Validation:} K-fold, stratified, bootstrapping
			\item \textcolor{myblue}{Grid Search:} Hyperparameter tuning
		\end{itemize}
		
		\vspace{0.1cm}
		\textcolor{myred}{\textbf{Exam Success:}} Understand methods, tradeoffs, and when to use each!
	\end{frame}
	
\end{document}